\documentclass[reqno,12pt]{article}

\usepackage[utf8]{inputenc}
\usepackage[tbtags]{amsmath}
\usepackage[english]{babel}
\usepackage{bm}
\usepackage{bbm} % \mathbbm{1}
\usepackage{amsfonts}
\usepackage{microtype}
\usepackage{physics}
\usepackage{tensor}
\usepackage{slashed}
\usepackage{subfig}
\usepackage{mathtools}
\usepackage[subfigure]{tocloft}
\usepackage[parfill]{parskip}
\usepackage{multirow}
\usepackage{xcolor}
\usepackage{biblatex}
\usepackage{verbatim} % for multiline comments
\usepackage{csquotes} % required by biblatex
\usepackage{graphicx}
\usepackage[font=small,labelfont=bf]{caption}
\usepackage{float}
\usepackage[toc,page]{appendix}
\usepackage{placeins}

\graphicspath{ {../pics/} }
\addbibresource{source.bib}

\numberwithin{equation}{section}

\renewcommand{\cftsecleader}{\cftdotfill{\cftdotsep}}

\newcommand{\red}[1]{\textbf{\textcolor{red}{#1}}}

\newcommand{\D}[1]{\,\mathcal{D}#1\,}
\newcommand{\Z}{\mathcal{Z}}
\newcommand{\R}{\mathbb{R}}
\newcommand{\SU}{\mathrm{SU}}
\newcommand{\Sp}{\mathrm{Sp}}
\newcommand{\id}{\mathbbm{1}}

\title{Tesi}
\date{2022}
\author{Nicholas Pini}

\begin{document}

\begin{titlepage}
	\begin{center}
		{{\Large{\textsc{Università di  Milano Bicocca}}}} \rule[0.1cm]{14cm}{0.1mm}
		\rule[0.5cm]{14cm}{0.6mm}
		\includegraphics[scale=0.45]{logo.png}\\
		{\small{\bf Dipartimento di Fisica G. Occhialini}}
					
	\end{center}
	\vspace{10mm}
	\begin{center}
			{\LARGE{Deconfinement transition in 3D Yang-Mills theory with Sp(2) gauge group and study of string effects}}\\
			\vspace{13mm}{\large Tesi di laurea magistrale}
	\end{center}
	\vspace{15mm}
	\par
	\noindent
	\begin{minipage}[t]{0.60\textwidth}
	{
		Relatore: prof. {\bf Leonardo Giusti}\\
		Correlatore: prof. {\bf Michele Pepe}\\
	}
	\end{minipage}
	\hfill
	\begin{minipage}[t]{0.40\textwidth}\raggedleft
			{Candidato: {\bf Nicholas Pini}\\
			Matricola: {\bf 813484}}
	\end{minipage}
	%\vspace{25mm}
	\begin{center}
	{Anno Accademico 2021/2022}
	\end{center}
\end{titlepage}

\begin{comment}
Le teorie di gauge non abeliane presentano il fenomeno del confinamento di colore: il potenziale d'interazione fra due cariche di colore cresce linearmente nel limite di lunghe distanze. Il tubo di flusso formato dalle linee di forza del potenziale può essere modellato come una stringa vibrante: questo modello è detto Effective String Theory (EST). È un modello molto predittivo a basse temperature e lunghe distanze, ma che perde affidabilità ad alte temperature, vicino al punto critico della transizione di deconfinamento. Nel regime critico, una congettura di Svetisky e Yaffe prevede che in una teoria di gauge (2+1) dimensionale, con transizione di deconfinamento di secondo ordine, i dettagli fini delle interazioni possono essere ignorati e il sistema è nella stessa classe di universalità del modello di Ising 2D, in quanto entrambi hanno la stessa simmetria Z2 e dimensionalità. Si è scritto un programma per simulare la teoria di Yang-Mills in tre dimensioni a temperatura finita regolarizzata su reticolo usando l'algoritmo Heat-Bath di Cabibbo e Marinari per verificare la congettura nel caso del gruppo di gauge Sp(2). Si è misurato il correlatore del loop di Polyakov (il parametro d'ordine della transizione di deconfinamento) e si è verificato che ha lo stesso andamento del correlatore fra spin del modello di Ising 2D, a corte e lunghe distanze, suggerendo che la congettura è vera per la teoria (2+1) dimensionale di Yang-Mills con gruppo di gauge Sp(2). 
\end{comment}

\newpage
\thispagestyle{empty}

\

\newpage

\begin{abstract}
Non abelian gauge theories exhibit color confinement: the potential between two color charges is linearly
rising in the large distance limit. An Effective String Theory (EST) exists which describes the flux tube formed by
the field lines of the potential as a vibrating string. This is an effective model with very good predictive power at
low temperatures and large distances, but it is not reliable in the high temperature regime, near the deconfinement critical point. Here,
the Svetisky-Yaffe conjecture predicts that a $(2+1)$ dimensional gauge theory is in the same universality class
of the 2D Ising model if the deconfinement phase transition is second order, 
because the fine details of the interactions can be neglected and the two systems share the same
$\mathbb{Z}_2$ symmetry and dimensionality. 
We wrote a computer simulation of the finite temperature 3D Yang-Mills theory regularized on a lattice,
in order to test this conjecture in the case
of the gauge group $\Sp(2)$, implementing the Heat-Bath algorithm by Cabibbo and Marinari. We measured the correlator
of Polyakov loops (the order parameter of the deconfinement phase transition) and verified that it behaves like the
spin-spin correlator of the 2D Ising model for short and large distances,
hinting that the conjecture is true for the $(2+1)$ Yang-Mills
gauge theory with gauge group $\Sp(2)$.
\end{abstract}

\newpage

\thispagestyle{empty}
\

\newpage

\tableofcontents

\newpage
\thispagestyle{empty}
\

\newpage


\section{Introduction}

Yang-Mills gauge theories exhibit color confinement: this is a phenomenon for which two color charges cannot be
observed separately because they interact through a confining interquark potential, which is linearly rising for
large distances. Despite a large body of numerical evidence, color confinement has no analytic proof.

Color confinement is a non perturbative phenomenon, so the natural context to study it is Lattice Gauge Theory (LGT) which is
a powerful regularization tool that discretizes spacetime in order to put a cutoff for the momentum of the theory. Together with
the path integral formulation of the theory, Wick rotated to be real and well defined, LGT is very suited to be simulated
on computers, which enables us to measure observables and better understand the physics behind color confinement. 

We have some insights about the interquark potential: a confining potential makes the Wilson loop, an observable defined
as a closed loop on the lattice, to follow an "area law", where the area is the one enclosed by said loop. In the
finite temperature LGT the role of the Wilson loop is taken by the correlator of the Polyakov loop, a closed loop that winds
around the periodic time direction of the lattice. It is in the finite temperature theory that we talk about deconfinement
phase transition: we find that the interquark potential is confining when the temperature of the system is below a certain
critical temperature; above this temperature, the system is instead in the deconfined phase, in which the potential is 
no longer confining. The order parameter of the deconfinement phase transition is the Polyakov loop, which is also 
associated to the spontaneous symmetry breaking of the center of the gauge group. An interesting conjecture by
Svetisky and Yaffe states that if the deconfinement transition is second order, then the $(d+1)$ dimensional
Yang-Mills gauge theory belongs to the same universality class of the $d$ dimensional spin model, with symmetry
the center of the gauge group. 

An effective approach to studying the interquark potential is the Effective String Theory (EST). The field lines of the 
potential form a flux tube, which the EST models as a vibrating string. It is only an effective model valid at low temperatures
and large distances, but it has had an incredible predictive power. This is thanks to the fact that the coefficients of the
first terms in the large distance expansion of this model are highly constrained by Poincaré invariance and match the one from
the Nambu-Goto model, the simplest EST, up to order $\flatfrac{1}{R^7}$ in the high temperature regime.
Despite its power, the EST is not consistent
with the picture by Svetisky and Yaffe: near the critical point of the deconfinement phase transition, the EST is no longer
as reliable. 

\cite{caselle} and \cite{caristo} studied this very topic in the context of the Yang-Mills theory defined for $\SU(N)$
gauge groups. Our goal is to instead tackle the problem using the $\Sp(N)$ gauge group, which is a subgroup of $\SU(2N)$ that
has the peculiarity of having the same center $\mathbb{Z}_2$ regardless of the size $N$ of the gauge group. In this study, 
we will focus on the case with $N=2$. 

This work is organized as follows: section \ref{background} lays the theoretical foundation for the study of the 
deconfinement phase transition (LGT, center symmetry, second order phase transitions and the Svetisky-Yaffe conjecture,
EST and a brief description of the $\Sp(2)$ group), section \ref{simulation} contains a detailed description of the
algorithm used to simulate the LGT, section \ref{results} is about the actual simulations that have been run and the data that
has been collected, and finally section \ref{conclusions} contains a recollection of the results and some final remarks.

\newpage

\section{Theoretical background} \label{background}

\subsection{Yang-Mills and Lattice gauge theory} \label{lgt}

Let us start by defining \textit{non abelian groups}. A group $G$ is said to be non abelian if there exists at least
a pair of group elements $g$ and $h$ belonging to $G$ for which

\begin{equation}
	[g, h] = gh - hg \ne 0.
\end{equation}

Yang-Mills theories are defined on non abelian unitary groups, called $\SU(N)$. 
$\SU(N)$ are compact, real Lie groups connected to the identity: these are groups of $N \times N$ matrices $U$
such that

\begin{equation}
	U^\dagger U = \id \qq{and} \det U = 1.
\end{equation}

In the fundamental representation, we consider $U$ to be parametrized with $\theta_a$ parameters, 
with $a = 1 \dots N^2 - 1$, and we define the generators of the Lie algebra to be

\begin{equation}
	iT_a \equiv \eval{\pdv{U}{\theta_a}}_{\theta_i = 0\,\forall i} \implies U = 1 + i \theta_aT^a,
\end{equation}

since $U$ is connected to the identity. Using the fact that $U^\dagger U = \id$, we have

\begin{equation}
	T^a = T^{a \dagger}.
\end{equation}

Thus, the generators are Hermitian. Near the identity we also have

\begin{equation}
	U = e^{i\theta_a T^a} \implies \det U = e^{i\theta_a \Tr T^a} = 1 \implies \Tr T^a = 0.
\end{equation}

Finally, $T^a$ are normalized such that

\begin{equation}
	\Tr(T^a T^b) = \frac{1}{2}\delta_{ab}.
\end{equation}

Note that the generators satisfy the commutation relation

\begin{equation}
	\qty[T^a, T^b] = if_{abc} T^c,	
\end{equation}

where $f_{abc}$ are the structure constants (which are totally antisymmetric) and depend on the algebra. 

The construction of the Yang-Mills action is based on gauge invariance. Consider a set of real functions $\Lambda^a(x)$ 
for $a = 1, \dots, N^2 - 1$, from which we can construct an element of the algebra doing a linear combination of generators
$\Lambda^a(x)T^a$. Exponentiating this element, we get 

\begin{equation} \label{eq:gauge_trans}
	G(x) = e^{i \Lambda_a(x) T^a} \in \SU(N),
\end{equation}

which belongs to the gauge group $\SU(N)$. Now, let us define the gauge field of the theory: 

\begin{equation}
	A_\mu(x) = A_\mu^a(x) T^a,
\end{equation}

where $A_\mu^a(x)$ is a real vector field. Following gauge invariance, we require the Yang-Mills action to be invariant
under the transformation

\begin{equation} \label{eq:A_mu_gauge_trans}
	A_\mu(x) \rightarrow G(x) A_\mu(x) G^\dagger(x) + i G(x) \partial_\mu G^\dagger(x).
\end{equation}

We also define

\begin{equation}
	F_{\mu\nu} = F_{\mu\nu}^a T^a \qq{where} 
	F_{\mu\nu} = \partial_\mu A_\nu - \partial_\nu A_\mu + i \qty[A_\mu, A_\nu].
\end{equation}

Note that $F_{\mu\nu}$ transforms under gauge transformations as

\begin{equation}
	F_{\mu\nu} \rightarrow G(x) F_{\mu\nu} G(x)^\dagger.
\end{equation}

The gauge invariant $D$ dimensional Yang-Mills action 
defined for non abelian unitary gauge groups is

\begin{equation} \label{eq:ym_action_mink}
	S_M = -\frac{1}{2g^2} \int \dd[D]{x_M} \Tr[F_{\mu\nu}F^{\mu\nu}].
\end{equation}

$g^2$ is the bare gauge coupling of the theory. The subscript "M" means that the action and the coordinates $x_M^\mu$ are defined
in Minkowski space time, whose line element is

\begin{equation} \label{eq:line_element}
	ds_M^2 = \qty(dx_M^0)^2 - \sum_{i=1}^D \qty(dx_M^i)^2.
\end{equation}

Let us now apply a \textit{Wick rotation}. This simply means defining new coordinates:

\begin{equation} \label{eq:wick_rotation}
	x_0^E = i x^0_M \qc x_i^E = x_M^i.
\end{equation}

This causes the line element in \eqref{eq:line_element} to become

\begin{equation}
	ds_E^2 = -\qty(dx^E_0)^2 - \sum_{i=1}^D \qty(dx^E_i)^2,
\end{equation}

which is the metric of the Euclidean space. Under this rotation, the derivates change to

\begin{equation}
	\partial_0^E = -i \partial_0^M \qc \partial_i^E = \partial_i^M.
\end{equation}

To keep the gauge transformation in \eqref{eq:A_mu_gauge_trans} as is, we also have to change the gauge field $A_\mu$:

\begin{equation}
	A_0^E = -i A^0 \qc A_i^E = -A^i_M.
\end{equation}

The Wick rotation has the effect of transforming the action in \eqref{eq:ym_action_mink} in a Euclidean version $S_E$:

\begin{equation}
	S_E = -i S_M.
\end{equation}

Thus, the Yang-Mills action in Euclidean space is

\begin{equation} \label{eq:ym_action}
	S_E = \frac{1}{2g^2} \int \dd[D]{x^E} \Tr[F_{\mu\nu}F_{\mu\nu}].
\end{equation}

with 

\begin{equation}
	\dd[D]{x^E} = i \dd[D]{x^M}.
\end{equation}

Ultimately, we are interested in measuring some observables that describe the Yang-Mills gauge theory, and we
do this in the path integral formulation. The idea is to
sum over all possible fields configuration weighted by the factor $e^{iS_M[A]} = e^{-S_E[A]}$.
This operation defines the \textit{path integral}:

\begin{equation} \label{eq:path_integral_mink}
	\Z = \int \D{A} e^{iS_M[A]} \qc \D{A} = \prod_{x, \mu} \delta A_\mu(x),
\end{equation}

with $\D{A}$ being gauge invariant. 
Under this formulation, we can calculate the expectation value of
some observable $\mathcal{O}$ as 

\begin{equation}
	\expval{\mathcal{O}} = \frac{1}{\Z} \int \D{A} e^{iS_M[A]} \mathcal{O}.
\end{equation}

Note that applying a Wick rotation, as defined in \eqref{eq:wick_rotation}, to
the path integral in \eqref{eq:path_integral_mink} transforms the exponential term $iS_M$ to $-S_E$,
a real quantity. The final form of the path integral we will be using is 

\begin{equation} \label{eq:path_integral}
	\Z = \int \D{A} e^{-S_E[A]} \qc \D{A} = \prod_{x, \mu} \delta A_\mu(x),
\end{equation}

and observables are calculated as

\begin{equation} \label{eq:obs_expval}
	\expval{\mathcal{O}} = \frac{1}{\Z} \int \D{A} e^{-S_E[A]} \mathcal{O}.
\end{equation}

We now introduce Lattice Gauge Theory (LGT), which is a way to regularize the theory. The basic concept
is to discretize spacetime with lattice spacing $a$ in a finite volume $V$, which yields a natural cutoff for the momentum:

\begin{equation}
	p \in \qty(-\frac{\pi}{a}, \frac{\pi}{a}),
\end{equation}

The idea is that, once the integral in \eqref{eq:path_integral} has been calculated, one rotates back to the original
theory and takes the continuum limit $a \rightarrow 0$ and $V \rightarrow \infty$. This regularization is particularly useful
for computer simulations, since we can now consider the theory to be a statistical system with Boltzmann factor of $e^{-S_E[A]}$
and use statistical methods, like Monte Carlo methods, to simulate the theory and compute observables. 

We need to find a way to define the action on a discretized lattice with $a \ne 0$ and finite volume $V$, 
in a way that we obtain the action in \eqref{eq:ym_action} in the continuum limit. Note that from now on we will drop the
superscript "E" and assume all quantities to be defined in the Euclidean space. 

On the lattice, $x^\mu$ now represents the position of a site, with $\mu = 0, \dots, d$ for a $(d+1)$ dimensional lattice.
We call $\hat{\mu}$ the unit vector in direction $\mu$: thus $x + \hat{\mu}$ is the site next to $x$ in the $\mu$ direction.  

Defining the discretized version of derivatives as

\begin{equation} \label{eq:discrete_derivative}
	\grad_\mu f(x) \equiv \frac{f(x + \hat\mu) - f(x)}{a},
\end{equation}

the discretized Yang-Mills action in \eqref{eq:ym_action} becomes

\begin{equation} \label{eq:discrete_ym_action}
	\widetilde{S} = \frac{1}{2g^2} a^4 \sum_x \sum_{\mu,\nu} \qty(\Tr\qty[F_{\mu\nu}F_{\mu\nu}] + O(a)),
\end{equation}

where $F_{\mu\nu}$ is now discretized as well:

\begin{equation} \label{eq:discrete_Fmunu}
	F_{\mu\nu}(x) = \grad_\mu A_\nu(x) - \grad_\nu A_\mu(x) - i\qty[A_\mu(x), A_\nu(x)].
\end{equation}

Noticing that $a^4 \sum_x$ is just the discretized version of the integral over spacetime, we can see that for
$a \rightarrow 0$, \eqref{eq:discrete_ym_action} becomes the original Euclidean Yang-Mills action in the continuum limit.
The problem with this approach is that $\widetilde{S}$ is gauge invariant \textit{only} in the limit $a \rightarrow 0$. For
$a \ne 0$, $\widetilde{S}$ is not gauge invariant. We want to change the discretization error in $O(a)$ in order to keep the 
gauge invariance even when $a \ne 0$.

To do so, let us introduce the \textit{link variable}:

\begin{equation} \label{eq:link}
	U_\mu(x) = e^{-iaA_\mu(x)} \in \SU(N) \text{ (the theory's gauge group)}.
\end{equation}

Being the exponentiation of an element of the algebra, the link variable belongs instead to the gauge group itself.
Note that the link is oriented: it starts at $x$, pointing along direction $\mu$, and connects to $x + \hat{\mu}$. Taking
the adjoint of a link variable is the same as considering the inverted link variable that goes from $x + \hat{\mu}$ to $x$:

\begin{equation}
	U_\mu(x)^\dagger = U_{-\mu}(x + \hat{\mu}).
\end{equation}

Links variable gauge transform as

\begin{equation} \label{eq:link_gauge_trans}
	U_\mu(x) \rightarrow G(x) U_\mu(x) G^\dagger(x + \hat{\mu}).
\end{equation}

Link variables are not gauge invariant. But we can construct invariant quantities under the gauge transformation in \eqref{eq:link_gauge_trans}
by considering closed loops of
ordered products of link variables, in which all terms $G(x)$ in \eqref{eq:link_gauge_trans} cancel out except the first and
last ones, and taking its trace. The cyclic property of the trace lets us cancel out the remaining $G(x)$ terms in the product,
leaving behind a gauge invariant quantity. With that said, we define the product of link variables around the smallest
loop possible, which is a square formed by four links:

\begin{equation} \label{eq:plaquette}
	U_{\mu\nu}(x) = U_\mu(x) U_\nu(x + \hat{\mu}) U_\mu^\dagger(x + \hat{\nu}) U_\nu^\dagger(x).
\end{equation}

This quantity is called \textit{plaquette}. We can see that under gauge transformations this quantity transforms as

\begin{equation} \label{eq:plaq_gauge_trans}
	\begin{aligned}
		U_{\mu\nu}(x) \rightarrow &G(x) U_\mu(x) G^\dagger(x + \hat{\mu}) \\ 
			&\times G(x + \hat{\mu}) U_\nu(x + \hat{\mu}) G^\dagger(x + \hat{\mu} + \hat{\nu})  \\
			&\times G(x + \hat{\mu} + \hat{\nu}) U_\mu^\dagger(x + \hat{\nu}) G^\dagger(x + \hat{\nu})  \\
			&\times G(x + \hat{\nu}) U_\nu^\dagger(x) G^\dagger(x) \\
			&= G(x) U_{\mu\nu}(x) G^\dagger(x),
	\end{aligned}	
\end{equation}

implying that the plaquette is gauge covariant. We can also see, as we just mentioned, that taking the trace of $U_{\mu\nu}$
cancels $G(x)$ and $G(x)^\dagger$ in \eqref{eq:plaq_gauge_trans}: thus, $\Tr U_{\mu\nu}(x)$ is a gauge invariant quantity.
We can now introduce the Wilson action:

\begin{equation} \label{eq:wilson_action}
	S_W = \frac{\beta}{N} \sum_x \sum_{\mu<\nu} \Re \Tr(\id - U_{\mu\nu}(x)),
\end{equation}

where $N$ is a normalization constant (the number of colors in the theory, related to the gauge group $\SU(N)$) and

\begin{equation} \label{eq:beta}
	\beta \equiv \frac{2N}{g^2},
\end{equation}

$g^2$ being the bare gauge coupling of the theory, the same as in \eqref{eq:ym_action}. \eqref{eq:wilson_action} defines
a gauge invariant action for lattice regularization at any value of $a \ne 0$, and reproduces \eqref{eq:ym_action} in the 
continuum limit.

Since $S_W$ involves only traces of plaquettes, it is clearly gauge invariant. We now show that it also reproduces the
Yang-Mills action for $a \rightarrow 0$. Consider the Campbell formula

\begin{equation}
	e^A e^B = \exp(A + B + \frac{1}{2}\qty[A, B] + \dots).
\end{equation}

Applying this formula iteratively, we write \eqref{eq:plaquette} as 

\begin{equation} \label{eq:expandend_plaquette}
	\begin{aligned}
		U_{\mu\nu}(x) &= \exp \biggl\{ -ia\biggl( A_\mu(x) + A_\nu(x + \hat{\mu}) - \frac{ia}{2}\qty[A_\mu(x), A_\nu(x+\hat\mu)] \\
		&- A_\mu(x + \hat\nu) - A_\nu(x) - \frac{ia}{2}\qty[A_\mu(x + \hat\nu), A_\nu(x)] \\
		&+ \frac{ia}{2}\qty[A_\mu(x) + A_\nu(x+\hat\mu), A_\mu(x + \hat\nu) + A_\nu(x)] + O(a^2) \biggr) \biggr\}.
	\end{aligned}
\end{equation}

If we do a Taylor expansion in $a$ for the gauge fields,

\begin{equation}
	A_\nu(x + \hat\mu) = A_\nu(x) + a \partial_\mu A_\nu(x) + O(a),
\end{equation}

we can keep only terms up to $a^2$ in \eqref{eq:expandend_plaquette}, which can then be rewritten as

\begin{equation} \label{eq:plaquette_in_Fmunu}
	\begin{aligned}
		U_{\mu\nu}(x) &= \exp{-ia^2\biggl( \grad_\mu A_\nu(x) - \grad_\nu A_\mu(x) - i\qty[A_\mu(x), A_\nu(x)] + O(a) \biggr)} \\
		&= \exp{-ia^2 \biggl(F_{\mu\nu}(x) + O(a)\biggr)},
	\end{aligned}
\end{equation}

recalling \eqref{eq:discrete_derivative} and \eqref{eq:discrete_Fmunu}. If we know substitute \eqref{eq:plaquette_in_Fmunu}
in $S_W$ and expand in $a$, we see that the constant terms cancel out, and the term of order $a^2$ is imaginary and thus removed by
the real part. Keeping only the term of order $a^4$, leaves us exactly with the discretized Yang-Mills action in 
\eqref{eq:discrete_ym_action}, which, as we have seen, reproduces the continuum theory in the limit $a \rightarrow 0$.

The theory is now suitable for being simulated on a computer: the Wick rotation guarantees that we are dealing with
real quantities, and the path integral formulation along with the gauge invariant Wilson action makes it possible
to consider the theory as a statistical system where fields are distributed with probability $e^{-S_W[U]}$.

It is important to note that the Wilson action in \eqref{eq:wilson_action} is not the only valid action for LGT, and
there are many ways to apply the lattice regularization to the gauge theory. Regardless, in the continuum limit
$a \rightarrow 0$, it must be the case that physical observables agree with measured observables: this implies that,
in general, bare quantities (like the coupling $g$) have a non trivial dependence on $a$, such that observables are 
independent of $a$ in the continuum limit (so that we retrieve the actual physical observables) and are finite. In other words
given a measured observable $\mathcal{O}$ and its actual physical value $\mathcal{O}_\text{phys}$, we have:

\begin{equation}
	\lim_{a \rightarrow 0} \mathcal{O}\qty(g(a), a) = \mathcal{O}_\text{phys},
\end{equation}

where $g(a)$ is the bare gauge coupling, now depending on $a$, in a way that keeps physics the same as we change $a$. In principle,
the form of $g(a)$ will depend on the particular observable $\mathcal{O}$ taken into consideration. However, for small enough values
of $a$, $g(a)$ should be universal, in order to ensure that physical observables are finite. This constraint on observables
can be formulated as

\begin{equation}
	a \dv{\mathcal{O}\qty(g(a), a)}{a} = 0,
\end{equation}

which can be rewritten as

\begin{equation} \label{eq:renormalization_group}
	\qty(\pdv{\ln a} - \beta(g)\pdv{g}) \mathcal{O}\qty(g(a), a) = 0,
\end{equation}

where 

\begin{equation} \label{eq:beta_function}
	\beta(g) \equiv -a \pdv{g}{a} = -\pdv{g}{\ln a}.
\end{equation}

\eqref{eq:renormalization_group} is called \textit{renormalization group equation} and depends on the $\beta$-function in
\eqref{eq:beta_function}, which is not the inverse coupling $\beta$ used in \eqref{eq:wilson_action}. The $\beta$-function
describes how the bare coupling $g$ depends on the lattice spacing $a$. For a pure gauge theory in $\SU(N)$, the $\beta$-function has
the form

\begin{equation}
	\beta(g) = -\beta_0 g^3 - \beta_1 g^5 + O\qty(g^7),
\end{equation}

where

\begin{equation}
	\begin{aligned}
		\beta_0 &= \frac{1}{16\pi^2} \frac{11}{3}N, \\
		\beta_1 &= \frac{1}{16\pi^2}\frac{34}{3}N^2
	\end{aligned}
\end{equation}

are universal coefficients which don't depend on the regularization scheme (this is not true for higher order terms). Using
\eqref{eq:beta_function}, we find

\begin{equation}
	a(g) = \frac{1}{\Lambda_L} \qty(\beta_0 g^2)^{\frac{\beta_1}{2\beta_0^2}} e^{-\frac{1}{2\beta_0 g^2}} \qty(1 + O\qty(g^2)),
\end{equation}

which can be inverted in order to find $g(a)$:

\begin{equation}
	g(a)^{-2} = \beta_0 \ln \qty(\frac{1}{a^2 \Lambda_L^2}) + \frac{\beta_1}{\beta_0} \ln \ln \qty(\frac{1}{a^2 \Lambda_L^2})
	+ O\qty(\frac{1}{\ln \qty(a^2 \Lambda_L^2)}).
\end{equation}

$\Lambda_L$ is an integration constant that sets the scale: the observables will have values in terms of this scale.
We can see that for $a \rightarrow 0$, we have $g(a) \rightarrow 0$. $g = 0$ is a zero for the $\beta$-function: thus, the
renormalization group equation is valid, and we have that the actual physics does not ultimately depend on $a$ in the continuum limit.
This behavior of the bare gauge coupling $g$ is called \textit{asymptotic freedom}, and it is typical of non abelian theories
for which the $\beta$-function is negative. Asymptotic freedom causes the gauge coupling to decrease as we decrease
the lattice spacing, meaning that we are making the lattice finer and finer. Note that lattice simulations can only
infer dimensionless ratios of physical quantities. This is because, in general, an observable $\mathcal{O}_\text{phys}$
will depend on $\Lambda_L$, as mentioned, in the following way:

\begin{equation}
	\mathcal{O}_\text{phys} = C \qty(\Lambda_L)^\delta,
\end{equation}

where $C$ is a constant and $\delta$ is the mass dimension of the observable. But recall that $\Lambda_L$ is not universal and
in general depends on the regularization scheme. Ratios of measured quantities, however, will not depend on $\Lambda_L$.

\subsection{Interquark potential} \label{potential}

Having now defined the LGT, we will be calling $N_s$ the length of every space dimension of the lattice (they 
all share the same length for simplicity's sake) and $N_t$ the length of the time dimension. With these definitions,
a $(d+1)$ dimensional lattice has a volume $V = (aN_t)(aN_s)^d$, where $a$ is the lattice spacing.

Consider two static color sources in the fundamental representation
on the lattice: to study the interquark potential that forms between them,
let us turn to \textit{Wilson loops}. A Wilson loop is defined as the trace of an ordered product 
of links along a closed path $\mathcal{C}$:

\begin{equation}
	W(\mathcal{C}) = \Tr[\prod_{(\mu,x)\in\mathcal{C}} U_\mu(x)].
\end{equation}

The loop $\mathcal{C}$ consists of four distinct pieces: two of them are so called \textit{Wilson lines}, which are 
products of links connecting two sites on the lattice sharing the same time coordinate; the other two pieces
are \textit{temporal transporters}: these are products of links forming a straight path connecting two sites at
different time coordinates, sharing the same spatial coordinates. If the Wilson lines are straight, the loop
$\mathcal{C}$ will form a square: this is called \textit{Planar Wilson loop}. Note that, being the trace of a closed
loop of links, the Wilson loop is gauge invariant.

To explain the physical interpretation of Wilson loops, let us briefly consider the theory to be in the temporal
gauge: this is a gauge in which all time-like links (and so the temporal transporters) are trivially the identity.
We can than write the expectation value of the Wilson loop, dropping the temporal transporters, as 

\begin{equation}
	\expval{W(\mathcal{C})} = \expval{W(\mathcal{C})}_\text{temp} = 
	\expval{\Tr[S(\va{x}, \va{y}, \flatfrac{\tau}{a})S(\va{x}, \va{y}, 0)^\dagger]}_\text{temp}
\end{equation}

where $S(\va{x}, \va{y}, \flatfrac{\tau}{a})$ is the Wilson line connecting the sites $x$ and $y$, which have the same time
coordinate $t = \frac{\tau}{a}$ n lattice units and spatial coordinates $\va{x}$, $\va{y}$, respectively. $\tau = at$ is then
the actual Euclidean time separation. The subscript "temp" is to
remind that we are using the temporal gauge (which doesn't change the expectations value of the theory). 
In Euclidean space, for large temporal extent $\tau$, we can write the correlator as

\begin{equation} \label{eq:wilson_lines_expansion}
	\expval{\Tr[S(\va{x}, \va{y}, \flatfrac{\tau}{a})S(\va{x}, \va{y}, 0)^\dagger]}_\text{temp} = 
	\sum_k \bra{0}\hat{S}(\va{x}, \va{y})\ket{k}\bra{k}\hat{S}(\va{x}, \va{y})^\dagger\ket{0} e^{-\tau E_k}.
\end{equation}

The states $\ket{k}$ describe a static quark-antiquark pair located at $\va{x}$ and $\va{y}$ \cite{gattringer}. 
Thus, the lowest energy state $E_1$ in \eqref{eq:wilson_lines_expansion} is the energy associated to this pair. 
Defining $R = a|\va{x} - \va{y}|$ to be the spatial distance between the two static quarks, we have

\begin{equation}
	E_1 = V(R).
\end{equation}

In other words, $E_1$ corresponds to the interquark potential. The expectation value of the Wilson loop
can be then written as

\begin{equation}
	\expval{W(\mathcal{C})} \sim e^{-\tau V(R)}.
\end{equation}

We can then retrieve the interquark potential as

\begin{equation}
	V(R) = - \lim_{\tau \rightarrow \infty} \frac{1}{\tau} \expval{W(\mathcal{C})}.
\end{equation}

The physical intuition behind Wilson loops is considering the static quark-antiquark pair as instantly
created at time $t = 0$ and moved apart at a distance $R$: the pair evolves with time until $t = \flatfrac{\tau}{a}$, at which
the pair is annihilated. 

In a confining gauge theory, the interquark potential is linearly rising for large values of $R$: 

\begin{equation}
	V(R) \sim \sigma_0 R,
\end{equation}

where $\sigma_0$ is a constant called \textit{zero temperature string tension}. Thus,

\begin{equation} \label{eq:area_law}
	\expval{W(\mathcal{C})} \sim e^{-\sigma_0 \tau R}.
\end{equation}

This is an "area law" for the expectation value of Wilson loop: the confining potential has as consequence
an exponent in \eqref{eq:area_law} of the form $\tau \times R$, which is the area of the Wilson loop enclosing
a rectangle on the lattice of with sides $\frac{\tau}{a}$, $\frac{R}{a}$ (using the number of links as unit). 

\subsection{Finite temperature LGT} \label{finitetemperature}

The description of the confining potential through Wilson loops does not take into account the temperature of the system.
If we want to study how the confinement-deconfinement transition depends on temperature, we need to define the theory at finite
temperature in the first place. 

It is well known how to do so for an LGT defined on a $(d+1)$ dimensional lattice of volume $(aN_s)^d(aN_t)$:
we consider the lattice regularization and impose \textit{periodic boundary condition} 
for bosonic fields (which are the fields we will be considering in this work) for the time dimension of the lattice. 
Now, the lattice represents a system of volume $(aN_s)^d$ and temperature

\begin{equation} \label{eq:lattice_temperature}
	T = \frac{1}{N_t a}.
\end{equation}

The compactified time direction is interpreted as the inverse of the temperature of the system.

In a finite temperature LGT, we cannot use Wilson loops anymore to describe the interquark potential. Now that the 
length of the time direction of the lattice is finite, we make the temporal extent of the Wilson loop as big as possible:
namely, we are setting $\tau = aN_t$. In this way, due to the periodic boundary condition,
the Wilson lines end up being on top of each other, oriented in opposite directions. We cannot use the temporal gauge
anymore, but we can gauge-transform the spatial pieces of the loop to the identity. The Wilson loop now reduces to two 
disconnected temporal transporters winding around the time direction at a spatial distance of $R$. We call the trace of one
of these temporal transporters \textit{Polyakov loop}:

\begin{equation} \label{eq:polyloop_definition}
	\phi(\va{x}) = \Tr[\prod_{j = 0}^{N_t - 1} U_0(j, \va{x})],
\end{equation}

where $U_0(j, \va{x})$ represents a time-like link at spatial coordinate $\va{x}$ and time coordinate $j$. The Polyakov
loop is a gauge invariant quantity.

Now, the expectation value of the Wilson loop is simply the expectation value of the product of the two Polyakov loops
$\phi(\va{x})$ and $\phi(\va{y})$ with $R = a |\va{x} - \va{y}|$, and as such, we find a relation for the interquark potential
$V(R, T)$ even in a finite temperature LGT:

\begin{equation} \label{eq:polyloop_correlator}
	\expval{\phi(\va{x}) \phi(\va{y})^\dagger} \sim e^{-aN_t V(R, T)} = e^{-\frac{1}{T} V(R, T)}.
\end{equation}

The interquark potential now depends also on the temperature $T$ of the system. Since we are considering a confining theory,
the interquark potential is linearly rising. We thus write, for large values of $R$,

\begin{equation}
	V(R, T) \sim \sigma(T) R,
\end{equation}

collecting the temperature dependence in $\sigma(T)$, the \textit{finite temperature string tension}. 
\eqref{eq:polyloop_correlator} becomes

\begin{equation}
	\expval{\phi(\va{x}) \phi(\va{y})^\dagger} \sim e^{-a\sigma(T)N_t R} = e^{-\frac{1}{T}\sigma(T) R}. 
\end{equation}

In the zero temperature limit $T \rightarrow 0$, we identify $V(R, T)$ with $V(R)$ and $\sigma(T)$ with $\sigma_0$.

The expectation value of a single Polyakov loop, $\expval{\phi(\va{x})}$, has an important physical role. In fact,
for large distances $R$, we expect the correlator of two Polyakov loops to factorize:

\begin{equation}
	\lim_{R \rightarrow \infty} \expval{\phi(\va{x})\phi(\va{y})^\dagger} = 
	\expval{\phi(\va{x})}\expval{\phi(\va{y})^\dagger} = |\expval{\phi}|^2.
\end{equation}

Due to translational invariance, the space position of $\phi$ in the last equality is not important, and we can simply
use the average of the Polyakov loop over the spatial volume of the lattice. Given \eqref{eq:polyloop_correlator}, for 
confining potentials, we have $\expval{\phi} = 0$, since $V(R, T) \rightarrow \infty$ with $R \rightarrow \infty$. Otherwise,
we have $\expval{\phi} \neq 0$. In other words, $\expval{\phi}$ is the \textit{order parameter} of the deconfinement phase
transition. Defining $F_q$ to be the free energy associated to a single color charge $q$, and interpreting the expectation
value of the Polyakov loop as the probability of observing a single static charge:

\begin{equation}
	\expval{\phi} \sim e^{-\frac{F_q}{T}},
\end{equation}

we can summarize the confinement-deconfinament transition as follows:

\begin{itemize}
	\item if $\expval{\phi} = 0$, the system is in the \textit{confined} phase, and $F_q \rightarrow \infty$;
	\item if $\expval{\phi} \neq 0$, the system is in the \textit{deconfined} phase, and $F_q$ is finite.
\end{itemize}

\subsection{Center symmetry} \label{center}

The Wilson action \eqref{eq:wilson_action} is not only gauge invariant, but it is equipped with a global symmetry, which
is realized through the center of the gauge group. Given a group $G$, the center $Z(G)$ is defined as

\begin{equation}
	Z(G) = \qty{z \in G \bigm| \forall g \in G \qc [z, g] = 0}.
\end{equation}

In other words, the center is the subgroup of $G$ which contains the set of all elements that commute with any element of the
entire group $G$. In the case of $\SU(N)$, the center of the group is $\mathbb{Z}_N$; for $Sp(N)$, the center is always
$\mathbb{Z}_2$, for all values of $N$. $\mathbb{Z}_N$ is a group defined as

\begin{equation}
	\mathbb{Z}_N = \qty{\exp(i \frac{2\pi k}{N}) \vdot \id_{N\times N} \bigm| k = 0, \dots, N - 1 }.
\end{equation}

The center transformation, under which the action is invariant, is defined as multiplying all time-like links in 
a fixed time slice $t = t_0$ by the same element $z$ belonging to the center of the gauge group:

\begin{equation}
	U_0(t_0, \va{x}) \rightarrow z U_0(t_0, \va{x}) \qc \forall \va{x} \qq{and fixed} t_0.
\end{equation}

Quantities built out of trivially closed loops of links, like Wilson loops and the Wilson action itself, are invariant under this 
transformation: this is because each loop has as many links going in one direction as links going in the opposite direction.
As such, the center transformation multiplies one time-like link pointing in one direction and the other time-link pointing
in the other direction and belonging to the same time slice. One of these two links is under the adjoint operation, and so
the center element $z$ multiplying this link is also adjoint. Since center elements commute with the entire group, $z$ and
$z^\dagger$ can be moved through the product of links in the loop and multiplied together. But $z z^\dagger = \id$, 
and thus the product of links is unchanged under center transformation.   

For example, consider the plaquette defined in \eqref{eq:plaquette} in the plane $(0, i)$, where $i > 0$ represents a spatial
direction, at the site $x = (t_0, \va{x})$. This would transform as

\begin{equation} \label{eq:center_trans}
	\begin{aligned}
	U_{0i}(x) &= U_0(t_0, \va{x}) U_i(t_0 + 1, \va{x}) U_0(t_0, \va{x} + \vu{i})^\dagger U_i(t_0, \va{x})^\dagger \\
	&\rightarrow z U_0(t_0, \va{x}) U_i(t_0 + 1, \va{x}) U_0(t_0, \va{x} + \vu{i})^\dagger z^\dagger U_i(t_0, \va{x})^\dagger \\
	&= z z^\dagger U_0(t_0, \va{x}) U_i(t_0 + 1, \va{x}) U_0(t_0, \va{x} + \vu{i})^\dagger U_i(t_0, \va{x})^\dagger \\
	&= U_{0i}(x).
	\end{aligned}
\end{equation}

However, the Polyakov loop is not invariant under center transformation. Looking at the way it is defined in
\eqref{eq:polyloop_definition}, we see that each link that appears in the product belongs to a different time slice: thanks
to the periodic boundary condition in the time direction, a Polyakov loop is closed without having links going in both direction,
as it is the case with Wilson loops and the plaquette. A Polyakov loop than transforms as

\begin{equation} \label{eq:polyloop_center_trans}
	\phi(\va{x}) \rightarrow z \phi(\va{x})
\end{equation}

Of course, the two point function of Polyakov loops in \eqref{eq:polyloop_correlator} is still invariant under center
transformations, since $z$ and $z^\dagger$ would cancel out. We can see that the center symmetry is spontaneously broken
when the system is above critical temperature and in the deconfined phase, because the Polyakov loop would have non zero 
expectation value. In the confined phase, below critical temperature, the expectation value of the Polyakov loop
is zero, which means the symmetry is unbroken. Therefore, the deconfinement phase transition signals a spontaneous symmetry
breaking of the center symmetry.  

\subsection{Second order phase transition} \label{phase_transition}

Now that we have discussed the role of the Polyakov loop and how it relates to confinement and the center symmetry, let us
discuss more broadly about second order phase transitions.

Phase transitions occur in systems in the thermodynamic limit: this is the limit in which the volume $V$ and
the number of particles $N$ go to infinity, while keeping $\flatfrac{V}{N}$ constant. However, we cannot calculate
observables nor simulate systems in the thermodynamic limit: we are forced to work in a finite volume, 
for a finite number of particles.
Thus, we instead study how some observables behave near the critical point (the temperature for which the system undergoes
the phase transition in the thermodynamic limit) at finite $V$ and $N$, and then extrapolate the behavior in the
thermodynamic limit. For example, second order phase transitions are characterized by a diverging susceptibility
and correlation length at the critical point, and by correlation functions which decay exponentially with distance. 

In our case, susceptibility is defined as

\begin{equation} \label{eq:susc_definition}
	\chi = \sum_{\va{x}} \expval{\phi(\va{0}) \phi(\va{x})} = N_s^d \expval{\phi^2}
\end{equation}

for the Polyakov loop $\phi$. The correlation length $\xi$ is a measure of the typical length at which parts of
the system interact with each other. If the deconfinement phase transition is second order, we then expect

\begin{equation}
	\chi \sim \qty(1 - \frac{T}{T_c})^{-\gamma} \qc \xi \sim \qty(1 - \frac{T}{T_c})^{-\nu},
\end{equation}

where $T$ is the temperature of the system and $T_c$ is the critical temperature.

The exponents with which these observables approach the critical point are called 
\textit{critical exponents}. $\nu$ and $\gamma$ are indeed example of critical exponents.
Additionally, second order phase transitions
are related to the spontaneous symmetry breaking of some symmetry of the system: typically, one considers an
order parameter of such symmetry breaking which reflects whether the symmetry has been broken. In our case,
as previously mentioned, the deconfinement
phase transition is related to the spontaneous symmetry breaking of the $\mathbb{Z}_2$ center symmetry, and the order 
parameter of such symmetry is the expectation value of the Polyakov loop $\phi$. When $\expval{\phi} = 0$, 
the center symmetry is unbroken, and the system is in the confined phase: this is the low temperature regime. When
the temperature is above $T_c$, $\expval{\phi} \ne 0$: the system is in the deconfined phase, in which the
center symmetry has been spontaneously broken. A critical exponent $\beta$ is also associated to the order parameter:

\begin{equation}
	\expval{\phi} \sim \qty(1 - \frac{T}{T_c})^\beta.
\end{equation}

In the deconfined phase, the order parameter $\phi$ thus settles on one of two values different from zero, due to
the $\mathbb{Z}_2$ broken symmetry. At least, this would happen in the thermodynamic limit. In actuality, tunneling events
can happen which causes the order parameter to jump from time to time between the two different values. This is
why in simulations we consider $\expval{|\phi|}$, the expectation value of the modulus of the Polyakov loop, to be the
order parameter: in a finite volume, $\expval{\phi} = 0$ even in the deconfined phase, while $\expval{|\phi|}$ is
non zero, as we would expect from the order parameter of the phase transition. 
The two values agree in the limit of infinite volume and are both zero in the confined phase. 

The importance of critical exponents lies in the concept of universality classes. Near the critical point, for systems
that have a second order phase transition, the fine details of the interactions can be neglected because the correlation
length $\xi$ diverges: the entire system interacts with itself, and small scale interactions are negligible. This causes
very different physical system, which have in common only their symmetry and dimensionality, to be in the same universality 
class and share the same critical exponents. We can then describe one system in the vicinity of the critical point using
functions calculated from another system belonging to the same universality class. For example, the 2D Ising model, which
describes the orientation of spin particles on a lattice, has a second order phase transition: from a disordered phase in
which spins are randomly oriented, the system passes to an ordered phase in which spins largely
line up in the same direction for distances of the order of $\xi$. This system also has a $\mathbb{Z}_2$ symmetry
given by the two possible values of the spin orientation, which is spontaneously broken in the in ordered phase. 
Its order parameter is the magnetization $m$, with which one also defines the susceptibility. This system has been
solved exactly, and its critical exponents are 

\begin{equation} \label{eq:critical_exponents}
	\nu = 1 \qc \gamma = \frac{7}{4} \qc \beta = \frac{1}{8}. 
\end{equation}

Note that in our case we are starting from a $(2+1)$ dimensional LGT, ending up with a 2D finite temperature gauge theory
with $\mathbb{Z}_2$ global symmetry: these are the same as the 2D Ising model. We will see that this correspondence 
is the core of the Svetisky and Yaffe conjecture.    

In order to define $T_c$, the critical temperature at which the system undergoes a phase transition, recall
the role that $\beta$ has in \eqref{eq:wilson_action}. For a fixed value of $N_t$, we call $\beta_c(N_t)$ the critical value
that $\beta$ has to have for the system to be in the critical point. We can then invert the function $\beta_c(N_t)$
and obtain $N_{t,c}(\beta)$: this is the critical value of the length of the time direction given a value of $\beta$.
Recalling how temperature is defined in a LGT \eqref{eq:lattice_temperature}, we can finally write an expression
for the critical temperature of the system:

\begin{equation} \label{eq:critical_temperature}
	T_c = \frac{1}{a N_{t,c}(\beta)}.
\end{equation}

Thus, given a value of $\beta$, we can determine at which temperature a phase transition occurs. This is an alternative way
to define $T_c$ that avoids changing the lattice spacing $a$, which will be set to $a = 1$ for ease of calculations. 


\subsection{The Svetitsky–Yaffe conjecture} \label{conjecture}

Given the discussion above about the role of Polyakov loops, it is interesting to study the confinement-deconfinement phase transition
using an effective action, built by integrating out the spacelike links of the LGT \cite{caristo} \cite{caselle}. This effectively
amounts to a \textit{dimensional reduction}: we end up with a $d$ dimensional spin model with global symmetry the center of the
gauge group, from a $(d+1)$ dimensional LGT, as first argued by Svetisky and Yaffe.

For a second order deconfinement transition, the long range dynamics is dominated by fluctuations in the Polyakov loop. Also,
interactions between Polyakov loops are suppressed exponentially in the long distance limit. Thus, in the vicinity of the 
critical point, the fine details of the Hamiltonian describing the effective spin model can be neglected, and the system will belong
to the same universality class of the simplest spin model, with the same spontaneous symmetry breaking and nearest neighbors interactions
only. For example, an SU(2) Yang-Mills theory in $(2+1)$, $(3+1)$ dimensions will be mapped to a $2$, $3$ dimensional spin model, with
global symmetry $\mathbb{Z}_2$, the center of the SU(2) gauge group. In particular, the 2 dimensional spin model is exactly integrable,
and it is an excellent starting pointing for studying the conjecture. If the conjecture holds, we expect
the critical exponents related to the deconfinement phase transition of a $(2+1)$ dimensional LGT to be the same as 
in \eqref{eq:critical_exponents}.

In particular, the conjecture claims the following:
\begin{itemize}
	\item The high temperature, deconfined phase of the original Yang-Mills theory corresponds to the low temperature, ordered phase of 
	the effective spin model. In both of these phases, the corresponding order parameter has non-zero expectation value.

	\item The correlator between Polyakov loops in the confining phase (which is the one we are interested in) is mapped to the spin-spin
	correlator in the disordered phase of the spin model. 
\end{itemize}

As the second point suggests, the spin-spin correlator will be an important quantity for testing the conjecture. Indeed, for the 2 dimensional
spin model these are well known (see \cite{caristo}). Defining $R$ to be the distance characterizing the two-point function, $\xi$ the
correlation length of the system and $\phi(x)$ the spin operator at $x$, we have
\begin{itemize}
	\item for $R \ll \xi$:
		\begin{equation} \label{eq:short_distance_spin}
			\begin{aligned}
				\expval{\phi(0) \phi(R)} &= \frac{k_s}{R^{\flatfrac{1}{4}}} \biggl[ 1 + \frac{R}{2\xi} \ln(\frac{e^{\gamma_e}R}{8\xi}) + \frac{R^2}{16\xi^2} \\
				&+ \frac{R^3}{32\xi^3} \ln(\frac{e^{\gamma_e}R}{8\xi}) + O\qty(\frac{R^4}{\xi^4}\ln^2\frac{R}{\xi}) \biggr]
			\end{aligned}
		\end{equation}

	\item for $R \gg \xi$:
		\begin{equation} 
			\expval{\phi(0) \phi(R)} = k_l K_0\qty(\frac{R}{\xi})
		\end{equation}
\end{itemize}

where $k_l$, $k_s$ are constants, $\gamma_e$ is the Euler-Mascheroni constant and $K_0$ is the modified Bessel function of order zero. 
Note that in the analysis of this work, we will be using, in the case $R \gg \xi$,
\begin{equation} \label{eq:large_distance_spin}
	\expval{\phi(0) \phi(R)} = k_l \qty(K_0\qty(\frac{R}{\xi}) - K_0\qty(\frac{N_s - R}{\xi})),
\end{equation}

where $N_s$ is the length of the space dimensions (which are all the same length). This is to account for the periodic boundary conditions
imposed on the simulated lattice.

\subsection{Effective string theory} \label{est}

A well established way to model two interacting quarks is through an Effective String Theory (EST). The field lines connecting
two static color charges in the fundamental representation of the gauge theory are concentrated in a flux tube: the basic
idea of the EST is to model this flux tube as a vibrating string, in the large distance and low temperature regime.

The first model for an EST was proposed by L{\"u}scher and collaborators: calling $(\xi^0, \xi^1)$ the coordinates
on the worldsheet and $X^i$, for $i = 1, \dots, D - 2$, the two dimensional bosonic fields, which describe the transverse
displacement of the worldsheet, they proposed the following Gaussian action:

\begin{equation} \label{eq:gaussian_action}
	S_\text{G}[X] = \frac{\sigma_0}{2} \int \dd[2]{\xi} \partial_a X^\mu \partial^a X_\mu + \text{perimeter term.}
\end{equation} 

Integrating this action leads to a correction term in the interquark potential, called L{\"u}scher term:

\begin{equation}
	V(R) \sim \sigma_0 R - \frac{\pi(D-2)}{24R},
\end{equation}

ignoring a constant term given by the perimeter term and terms of order beyond $\flatfrac{1}{R}$. The integration
can be done in the finite temperature case, also. In the low temperature regime, this yields only a minor correction
to the L{\"u}scher term; in the high temperature regime (but still below the deconfinement critical point), the interquark
potential assumes the form \cite{caselle}

\begin{equation}
	V(R, T) \sim \qty(\sigma_0 - \frac{\pi(D - 2)}{6N_t^2})R.
\end{equation}

The action in \eqref{eq:gaussian_action} is not consistent with the constraints imposed by Lorentz invariance. In fact,
it is actually the first term of the large distance expansion of the Nambu-Goto action:

\begin{equation} \label{eq:nambugoto_action}
	S_\text{NG} = \sigma_0 \int_\Sigma \dd[2]\xi \sqrt{g},
\end{equation}

where $g = \det g_{\alpha\beta}$ and $g_{\alpha\beta} = \partial_\alpha X_\mu \partial_\beta X^\mu$. $\Sigma$ is the surface of the
world-sheet, on which are defined the coordinates $\xi = (\xi^0, \xi^1)$. \eqref{eq:nambugoto_action} represents the 
simplest EST model fulfilling the Lorentz constraints, and it is simply the generalization of the relativistic
action for a point-like particle.
In fact, every possible configuration of the string is weighted proportionally to the area spanned by
the string. It is quite standard to use this model in the "physical gauge", in which 
\begin{equation}
	\xi^0 = X^0 \qc \xi^1 = X^1.
\end{equation} 

In other words, the worldsheet coordinates are identified with the longitudinal degrees of freedom of the string: the string action is left with
$(D-2)$ degrees of freedom representing the transverse displacement. 

Expanding the Nambu-Goto model in the low energy regime \cite{caselle}, shows that all other terms
beyond the Gaussian one combine to give an exactly integrable and irrelevant perturbation of the Gaussian action.
This interesting property makes it possible to exactly compute the partition function of the model, and to exactly write the two-point
function of the Polyakov loop in $D$ dimensions as

\begin{equation} \label{eq:polyloop_expansion}
	\expval{\phi(0) \phi(R)} = 
	\sum_{n=0}^{\infty} w_n \frac{2R\sigma_0N_t}{E_n} \qty(\frac{\pi}{\sigma_0})^{\frac{D-2}{2}} 
	\qty(\frac{E_n}{2\pi R})^{\frac{D-1}{2}} K_{\flatfrac{(D-3)}{2}}(E_n R),
\end{equation}

where $K_\nu$ is the modified Bessel function, $w_n$ are the weights (in the $D=3$ case, they are equal to the number of partitions
of the integer $n$), and $E_n$ are the energy levels:

\begin{equation}
	E_n = \sigma_0 N_t \sqrt{1 + \frac{8\pi}{\sigma_0 N_t^2}\qty(n - \frac{D-2}{24})}.
\end{equation}

At large distances $R$, the two points function is dominated by the lowest state

\begin{equation}
	E_0 = \sigma_0 N_t \sqrt{1 - \frac{\pi(D - 2)}{3\sigma_0 N_t^2}}
\end{equation}

and \eqref{eq:polyloop_expansion} goes like

\begin{equation} \label{eq:est_polyloop_corr}
	\expval{\phi(0) \phi(R)} \sim \qty(\frac{1}{R})^{\frac{D-3}{2}} K_{\flatfrac{(D-3)}{2}}(E_n R).
\end{equation}

Remarkably, this is exactly the large distance behavior of the spin-spin correlator in \eqref{eq:large_distance_spin}
with $D = 3$.
This is true in general: any spin model with an isolated ground state is described
by a modified Bessel function, and the general form of \eqref{eq:polyloop_expansion} is fixed by the EST description itself,
with very mild assumptions: only the weights change.
In this picture, $E_0$ is thus the inverse of the correlation length $\xi$. 

Setting $D = 3$ and defining 

\begin{equation}
	\sigma(T) = \sigma_0 \sqrt{1 - \frac{\pi}{3 \sigma_0 N_t^2}}, 
\end{equation}

we have

\begin{equation} \label{eq:E0_with_nt}
	E_0 = \frac{1}{\xi} = \sigma(T) N_t.
\end{equation}

We can extrapolate the critical exponent $\nu$ of the phase transition: this represents how the correlation length
$\xi$ approaches the critical point of the system. To do so, we write the critical temperature as

\begin{equation}
	T_c = \sqrt{\frac{3\sigma_0}{\pi}}.
\end{equation}

Using \eqref{eq:E0_with_nt} and the $T$-$N_t$ relation written in \eqref{eq:lattice_temperature}, we get

\begin{equation}
	E_0 = \frac{1}{\xi} = \sigma_0 N_t \sqrt{1 - \frac{T^2}{T_c^2}}.
\end{equation}

This implies

\begin{equation}
	\xi \sim \qty(1 - \frac{T^2}{T^2_c})^{-\flatfrac{1}{2}} = \qty(1 - \frac{T^2}{T^2_c})^{-\nu}.
\end{equation}

The correlation length $\xi$ approaches the critical point with a critical exponent $\nu = \flatfrac{1}{2}$. This signals
a second order phase transition, and it is typical of the mean field approximation. But this is in contrast with the
Svetisky and Yaffe analysis: a $(2+1)$ dimensional gauge theory ($D=3$), with a second order
deconfinement phase transition, should be in the same universality class of the 2D Ising model, which has a critical
exponent for the correlation length of $\nu = 1$. This means that the Nambu-Goto EST model is not completely reliable
in the vicinity of the critical point, despite being a very good approximation in the low temperature regime. 

In fact, as mentioned in \cite{caselle} and \cite{caristo}, corrections to the Nambu-Goto EST model, which are of
particular interested in modern physics, has been found by reasoning about the coefficients of the large distance expansion
and their constraints. It can be shown that the first correction to the Nambu-Goto model is of order $\flatfrac{1}{R^7}$: this
is why the Nambu-Goto model and its first order expansion Gaussian term have been quite successful in describing the
interquark potential at large distances.   

\subsection{Sp(2) group} \label{sp2}

In this work, we will focus on Yang-Mills LGT with the $\Sp(2)$ gauge group: not only this deviates from the usual choice of $\SU(N)$, but
it has the important property of having $\mathbb{Z}_2$ as center symmetry, regardless of the dimension of the group itself. Also,
it is known that the deconfinement phase transition of $\Sp(2)$ Yang-Mills gauge theory is second order. 

Let us lay out a brief description of the $\Sp(N)$ group. $\Sp(N)$ is a subgroup of $\SU(2N)$, in which its elements leave the skew-symmetric
matrix $J$ invariant, where

\begin{equation}
	J = \mqty(0 & \id_{N \times N} \\ -\id_{N \times N} & 0) = i\sigma^2 \otimes \id_{N \times N}.
\end{equation}

Note that $J$ also belongs to $\Sp(N)$.

Given $U \in \Sp(N)$ (which, of course, also implies $U \in \SU(2N)$), it holds that

\begin{equation}
	U^* = J U J^\dagger.
\end{equation}

$U$ and $U^*$ are related by a unitary transformation, and thus the 2N dimensional fundamental representation of $\Sp(N)$ is pseudo-real.
Also, charge conjugation is simply a global gauge transformation.

Let us show that $\Sp(N)$ is indeed a group. First of all, the identity matrix belongs to the group. Given $U \in \Sp(N)$, $U^\dagger$ also
belongs to the group:

\begin{equation}
	(U^\dagger)^* = (U^*)^\dagger = (J U J^\dagger)^\dagger = J U^\dagger J^\dagger \implies U^\dagger \in \Sp(N).
\end{equation}

Finally, a product of $\Sp(N)$ elements $U$, $V$ belongs to the group itself:

\begin{equation}
	(U V)^* = U^* V^* = J U J^\dagger J V J^\dagger = J U V J^\dagger \implies UV \in \Sp(N).
\end{equation}

Considering the constraints above, a generic element of $\Sp(N)$ has complex entries of the form

\begin{equation} \label{eq:sp2_matrix_form}
	U = \mqty(W_{N \times N} & X_{N \times N} \\ -X_{N \times N}^* & W_{N \times N}^*).
\end{equation}

Note, however, that $U$ also must belong to $\SU(2N)$, which implies that $U$'s eigenvalues come in conjugate pairs. Since center elements
are multiples of the identity matrix, we have that $W_{N \times N} = W_{N \times N}^*$ in \eqref{eq:sp2_matrix_form}, 
and thus the center of $\Sp(2)$ is $\mathbb{Z}_2$.

In the specific case of $\Sp(2)$, which is the gauge group we will be considering, the generic matrix form we
will be using is

\begin{equation}
	U = \mqty(
		W_{11} & W_{12} & X_{11} & X_{12} \\
		W_{21} & W_{22} & X_{21} & X_{22} \\
		X^*_{22} & -X^*_{21} & W^*_{22} & -W^*_{21} \\
		-X^*_{12} & X^*_{11} & -W^*_{12} & W^*_{11}
	),
\end{equation}

where each of its entries is a complex number. Finally, figure \ref{fig:sp2weights} shows the weight diagram of the
fundamental representation of the $\Sp(2)$ group. 

\begin{figure}[h]
	\centering
	\includegraphics[width=0.5\textwidth]{sp2_weights.png}
	\caption[$\Sp(2)$ weight diagram]{Weight diagram of the fundamental representation of the $\Sp(2)$ group.}
\label{fig:sp2weights}
\end{figure}

\newpage

\section{Simulation and algorithm} \label{simulation}

As explained in section \ref{lgt}, we define the partition function of the $(2+1)$ dimensional system to be

\begin{equation}
	\mathcal{Z} = \int \D{U} e^{-S_W[U]} \qq{where} \D U \equiv \prod_{x,\mu} \dd{U_{x,\mu}}.
\end{equation}

$U_{x,\mu}$ (a link variable starting at the lattice site $x$, pointing
along the direction $\mu$) is an element of the gauge group $\Sp(2)$ introduced in section \ref{sp2}
and $S$ is the action, defined as
\begin{equation} \label{eq:algorithm_action}
	S_W = - \frac{\beta}{4} \sum_\square S_\square \qc 
	S_\square \equiv \Tr(U_{x,\mu}U_{x+\hat{\mu}, \nu}U^\dagger_{x+\hat{\nu}, \mu}U^\dagger_{x,\nu}).
\end{equation}
The symbol $\square$ represents the plaquette: the minimum loop possible on the lattice. In the case of the $\Sp(2)$
gauge group ($\Sp(N)$ with $N = 2$), the trace normalization is $2N = 4$, which is why $\beta$ in \eqref{eq:algorithm_action} has
4 in the denominator. In fact, $\beta$ is now defined as

\begin{equation} \label{eq:algorithm_beta}
	\beta = \frac{8}{g^2}.
\end{equation}

As opposed to \eqref{eq:wilson_action}, in \eqref{eq:algorithm_action} we ignored the numerical factors deriving from
the trace of the identity. Also, the trace of $\Sp(2)$ group elements is always real, so we also removed the real part
from the expression of the action, because it would be redundant. 

Recall that a link 
variable has a direction: the adjoint of a link variable is the link connecting the two sites in the opposite
direction. In other words:
\begin{equation}
	U_\mu(x)^\dagger = U_{-\mu}(x + \hat\mu).
\end{equation}

\subsection{Monte Carlo methods} \label{montecarlo}

We have seen in section \ref{lgt} that we are able to formulate the Yang-Mills theory in a way
that makes the system treatable as a statistical system, where the fields are distributed according to the Boltzmann
factor $e^{-S_W[U]}$: we can then use Monte Carlo methods to solve the integrals in the theory. 

Recall that, given the path integral $\Z$ over some configuration of the links $U_\mu(x)$,

\begin{equation}
	\Z = \int \D U e^{-S_W[U]} \qc \D U = \prod_{x,\mu} \delta U_\mu(x),
\end{equation}

we can calculate the expectation value of an observable $\mathcal{O}$ as

\begin{equation} \label{eq:montecarlo_obs_expval}
	\expval{\mathcal{O}} = \frac{1}{\Z} \int \D U e^{-S_W[U]} \mathcal{O}.
\end{equation}

To exactly solve this integral, we would need to sum over all possible configurations of the lattice: this task is
clearly impossible. However, we can tackle the problem using Monte Carlo methods: 
these approximate solutions to integrals based on probability theory, without needing
to consider all possible configurations. We consider the links $U_\mu(x)$ to be random variables distributed as

\begin{equation} \label{eq:mc_links_probdist}
	\dd{P(U)} = \frac{1}{\Z} e^{-S_W[U]} \D U.
\end{equation}

The Monte Carlo method of \textit{importance sampling} estimates the integral in \eqref{eq:montecarlo_obs_expval} as

\begin{equation} \label{eq:mc_estimate_obs}
	\expval{\mathcal{O}} = \lim_{N \rightarrow \infty} \frac{1}{N} \sum_{n=1}^N \mathcal{O}(U).
\end{equation}

The idea behind importance sampling is to evaluate observables in randomly generated configurations, which are
distributed according to the Boltzmann factor in the path integral. Configuration weighted more heavily by $e^{-S_W[U]}$
are then more probable to show up and end in up in the average \eqref{eq:mc_estimate_obs}: thus, in the limit of
$N \rightarrow \infty$, we retrieve the observable as distributed according to \eqref{eq:montecarlo_obs_expval}. Of course,
in actual simulations $N$ must be finite: in this case, the uncertainty of $\expval{\mathcal{O}}$ behaves like 
$O\qty(\flatfrac{1}{\sqrt{N}})$. Importance sampling is extremely powerful: we simply need to simulate the lattice
in order to have links distributed as \eqref{eq:mc_links_probdist} and measure some observable $\mathcal{O}$.
Doing it enough times and averaging over all values of $\mathcal{O}$, will yield an accurate estimate for the observable.

We now face another problem, however. We cannot simply pick a random configuration from the space of all possible
configurations according to the probability distribution in \eqref{eq:mc_links_probdist}. But we can
generate a configuration starting from another, previous lattice configuration, using \textit{Markov chains}.

Markov chains involve a starting configuration of random variables on the
lattice - let us call it $L_0$ - from which subsequent configurations $L_i$ are generated in order, one after the other.
This process is constructed in a way that makes each possible configuration reachable by the Markov process, giving more
importance (and thus more probability of being obtained) to configurations with a larger weight $e^{-S_W[U]}$. The 
probability of getting a configuration $L_i$ only depends on the previous configuration $L_{i-1}$ of the lattice:
we will call this $Q(L_i|L_{i-1})$. These probabilities $Q(L_i|L_{i-1})$ are named \textit{transition probabilities}.
For a Markov process, it must hold

\begin{equation}
	0 \le Q(L'| L) \le 1 \qc \sum_{L'} Q(L'|L) = 1 \qc \forall L,L'. 
\end{equation}

This is just saying that any configuration $L$ has the ability to transition to any other configuration $L'$, including
the case $L = L'$.

Let $P(L)$ be the probability for the system to be in the configuration $L$. It is possible to show using the normalization
property above that

\begin{equation}
	P(L') = \sum_L Q(L' | L).
\end{equation}

It is important to note that we want the Markov process to be in equilibrium, in order to start measuring observables.
This is achieved when 

\begin{equation} \label{eq:balance_equation}
	\sum_L Q(L'| L) P(L) = \sum_L Q(L| L') P(L').
\end{equation}

We are essentially requesting that the probability of arriving at configuration $L'$ is the same as the probability
of leaving the configuration $L'$. If this is true, then $P(L)$ is a fixed point of the Markov process: once our system
reaches this stage by applying the transition probability $Q$ various times, it is in equilibrium and 
the probability $P(L)$ doesn't change by applying $Q$ again. This point of equilibrium is often called \textit{thermalization}.
There is not a way to precisely decide
when the lattice simulation is thermalized, and normally one observes some observable of interested and see after how many
steps it settles around some average value.

A simple solution to \eqref{eq:balance_equation} is

\begin{equation} \label{eq:detailed_balance_condition}
	Q(L'| L) P(L) = Q(L| L') P(L').
\end{equation}

This is called \textit{detailed balance condition}. This is a sufficient condition that assures the system will
reach an equilibrium at some point. As mentioned above, we also want every configuration to be reachable: this property
is called \textit{ergodicity} and it is obtained if $Q(L|L') > 0$ for any pair $L$, $L'$.

\subsection{Heat-bath algorithm} \label{heatbath}

In this section, we will describe in detail the algorithm used to generate subsequent lattice configurations respecting
the requirements of Markov processes, called \textit{Heat-Bath}. 
We will first introduce the algorithm as used for the $\SU(2)$ gauge group,
which has been designed by Creutz \cite{creutz}, and then we will describe 
how it translates to our gauge group of interest, $\Sp(2)$: this will be the algorithm used in our
simulation, and it has been designed by Cabibbo and Marinari \cite{cabibbo}.

Consider a LGT with gauge group $\SU(2)$, using the Wilson action as in \eqref{eq:wilson_action}. We call $u$
link variables that belongs to $\SU(2)$: these represent the random variables we want distributed as

\begin{equation} \label{eq:u_probdist}
	\dd{P(u)} = \frac{1}{\Z} e^{-S_W[u]} \dd{u}.
\end{equation}

If we focus on a single link that needs to be updated, we only need to consider the plaquettes in which it is
present: the rest of the action does not depend on this particular link. We call \textit{staple} the product
of the three links which, together with the link we are updating, forms a plaquette. In a 3D lattice, each link
is surrounded by 4 staples (we don't consider links on the perimeter of the lattice, because our simulations will 
use periodic boundary conditions for all dimensions). We can then write the action as

\begin{equation} \label{eq:heatbath_wilson_action_su2}
	S_W[u] = -\frac{\beta}{4} \Tr(u \sum_i \tilde{u}_i) + \text{terms independent of } u,
\end{equation}

where $\tilde{u}_i$ represents one staple, and $i = 1, \dots, 4$. Note that we are
using the normalization for the $\Sp(2)$ in \eqref{eq:heatbath_wilson_action_su2} despite being defined for
$\SU(2)$: the reason is to avoid confusion with the already defined $S_W$ and to be consistent with the actual implemented
algorithm. Apart from a numerical factor, it does not change Creutz's algorithm in any way. 

Ignoring normalization constants, we can rewrite \eqref{eq:u_probdist} as

\begin{equation} \label{eq:u_probdist_2}
	\dd{P(u)} \sim \exp\qty[\frac{\beta}{4} \Tr(u \sum_i \tilde{u}_i)] \dd{u}.
\end{equation}

Since $u \in \SU(2)$, we parametrize it in the following way:

\begin{equation} \label{eq:su2parametrization}
	u = \alpha_0 \id + i \va{\alpha} \cdot \va{\sigma},
\end{equation}

where $\alpha_\mu \in \mathbb{R} \ \forall \mu = 0, 1, 2, 3$ with the constraint that 

\begin{equation}
	\alpha^2 \equiv \alpha_0^2 + |\va{\alpha}|^2 = 1
\end{equation} 

and $\va{\sigma} = \qty(\sigma_1, \sigma_2, \sigma_3)$ is the three-vector of $2 \times 2$ Pauli matrices. 

The $\SU(2)$ group measure is then

\begin{equation} \label{eq:su2_group_measure}
	\dd{u} = \frac{1}{2\pi^2} \delta(\alpha^2 - 1) \dd[4]{\alpha}.
\end{equation}

Since the sum of $\SU(2)$ elements is proportional to an $\SU(2)$ element, we write

\begin{equation} \label{eq:su2_sum}
	\sum_i \tilde{u}_i = c \bar{u} \qc \bar{u} \in \SU(2)
\end{equation}

where 

\begin{equation}
	c = \det(u \sum_i \tilde{u}_i)^{\flatfrac{1}{2}}.
\end{equation}

Inserting \eqref{eq:su2_sum} in \eqref{eq:u_probdist_2}, we have

\begin{equation}
	\dd{P(u)} \sim e^{\frac{1}{4}\beta\Tr(c u \bar{u})} \dd{u}.
\end{equation}

The group measure in \eqref{eq:su2_group_measure} is invariant under multiplication by another $\SU(2)$ element:

\begin{equation}
	\dd(b u) = \dd{u} \qfor b \in \SU(2),
\end{equation}

so we can write

\begin{equation} \label{eq:probdist_u_ubar}
	\dd{P(u \bar{u}^{-1})} \sim e^{\frac{1}{4}\beta c \Tr(u)} \dd{u} 
	= \frac{1}{2\pi^2} e^{\frac{\beta}{2}c \alpha_0} \delta\qty(\alpha^2 - 1) \dd[4]{\alpha},
\end{equation}

because $\Tr(u) = 2\alpha_0$. Noticing that $\delta\qty(\alpha^2 - 1)\dd[4]{\alpha} = 
\frac{1}{2}\qty(1 - \alpha_0^2)^{\flatfrac{1}{2}} \dd{\alpha_0}\dd{\Omega}$, we rewrite \eqref{eq:probdist_u_ubar} as

\begin{equation}
	\dd{P(u \bar{u}^{-1})} \sim 
	\frac{1}{2\pi^2} \frac{1}{2}\qty(1 - \alpha_0^2)^{\flatfrac{1}{2}} e^{\frac{\beta}{2}c \alpha_0} \dd{\alpha_0}\dd{\Omega},
\end{equation}

where $\alpha_0 \in (-1,1)$ and $\dd{\Omega}$ is the differential solid angle of the three-vector $\va{\alpha}$,
which is of length $\qty(1 - \alpha_0^2)^{\flatfrac{1}{2}}$.

The problem is now about generating the four-vector $\alpha_\mu$ according to the distribution above.
We start by randomly generating $\alpha_0$ according to

\begin{equation} \label{eq:probdistalpha}
	P(\alpha_0) \sim (1 - \alpha_0^2)^{\flatfrac{1}{2}} e^{\frac{\beta}{2}c\alpha_0}.
\end{equation}

The algorithm is quite simple. We uniformly generate $x$ in the range 

\begin{equation}
	e^{-\beta c} < x < 1
\end{equation} 

and define a trial $\alpha_0$ distributed according to $e^{\frac{\beta}{2}c\alpha_0}$ as

\begin{equation}
	\alpha_0 = 1 + \frac{2}{\beta c}\ln{x}.
\end{equation}

To account for the term $\qty(1 - \alpha_0^2)^{\flatfrac{1}{2}}$ in \eqref{eq:probdistalpha}, we \textit{reject} this trial
$\alpha_0$ with probability $1 - \qty(1 - \alpha_0^2)^{\flatfrac{1}{2}}$, generating a new trial $\alpha_0$ if 
the rejection is successful. We keep doing this until a trial $\alpha_0$ is finally accepted. 

The unit vector $\va{\alpha} = (\alpha_1, \alpha_2, \alpha_3)$ is constructed by uniformly generating 

\begin{equation}
	\begin{aligned}
		\phi &\in (0, 2\pi), \\
		y &\in (-1, 1)
	\end{aligned}
\end{equation}

and defining 

\begin{equation}
	\begin{aligned}
		\theta &\equiv \arccos(y), \\
		r &\equiv \qty(1 - \alpha_0)^{\flatfrac{1}{2}},
	\end{aligned}
\end{equation}

$\alpha_0$ being the trial random number that has been accepted according to \eqref{eq:probdistalpha}. 

We finally have, for $\va{\alpha} = (\alpha_1, \alpha_2, \alpha_3)$,

\begin{equation}
	\begin{cases}
		\alpha_1 = r \sin(\theta) \cos(\phi) \\
		\alpha_2 = r \sin(\theta) \sin(\phi) \\
		\alpha_3 = r \cos(\theta)
	\end{cases}.
\end{equation}

Having reconstructed the full vector $\alpha_\mu$, we can get $u$ using \eqref{eq:su2parametrization}: the new link
variable on the lattice will finally be $u \bar{u}^{-1}$.

Creutz's algorithm is very efficient in generating new lattice configurations, but it relies heavily on the $\SU(2)$
property for which a sum of $\SU(2)$ is proportional to an $\SU(2)$ element. We will now explain how Cabibbo and Marinari
managed to extend Creutz's heat-bath algorithm to other gauge groups. In particular, we will apply it to $\Sp(2)$, the gauge
group of the Yang-Mills theory we want to study.

The idea behind Cabibbo and Marinari's Heat-Bath algorithm is to apply Creutz's algorithm to different $\SU(2)$
subgroups of the theory's gauge group. Of course, multiple subgroups are necessary, in order to ensure ergodicity.
Given a link to update, we extract an $\SU(2)$ element belonging to one of the subgroups, apply Creutz's algorithm,
embed the result again in $\Sp(2)$ and finally left multiply the original link. We iterate this procedure for each
subgroup, until the original link is completely updated. 


We consider a set $F$ of $\SU(2)$ subgroups of the gauge group $\Sp(2)$, with the requirement that there is no subset of $\Sp(2)$
that is invariant under left multiplication by $F$: this makes sure that the algorithm is ergodic.
Given a $\Sp(2)$ element $U$ in the fundamental representation of the form

\begin{equation}
	U = \mqty(
		W_{11} & W_{12} & X_{11} & X_{12} \\
		W_{21} & W_{22} & X_{21} & X_{22} \\
		X^*_{22} & -X^*_{21} & W^*_{22} & -W^*_{21} \\
		-X^*_{12} & X^*_{11} & -W^*_{12} & W^*_{11}
	),
\end{equation}

where $W_{ij}, X_{kl} \in \mathbb{C}$ for $i,j,k,l = 1, 2$, we have four ways of extracting two complex numbers $t_1$ and $t_2$:

\begin{itemize}
	\item $\begin{cases} t_1 = W_{11} \\ t_2 = X_{12} \end{cases}$
	\item $\begin{cases} t_1 = W_{22} \\ t_2 = X_{21} \end{cases}$
	\item $\begin{cases} t_1 = W_{11} + W_{22} \\ t_2 = X_{11} - X_{22} \end{cases}$
	\item $\begin{cases} t_1 = W_{11} + W^*_{22} \\ t_2 = W_{12} - W^*_{21} \end{cases}$
\end{itemize}

We can build an $\SU(2)$ group $a_k$ element as

\begin{equation} \label{eq:su2element}
	a_k = \mqty(t_1 & t_2 \\ -t_2^* & t_1^*) \qc k = 1, 2, 3, 4,
\end{equation}

where $k$ labels each subgroup. 
Each choice of $t_1$ and $t_2$ above gives a different $\SU(2)$ element belonging to a $\SU(2)$ subgroup of $\Sp(2)$.

These come from considering the weights of the representation of the group as shown in figure \ref{fig:sp2weights} and from
the fact that we ultimately only care about what ends up in the trace of $\Sp(2)$ elements. This choice of subgroup
guarantees that each part of the $\Sp(2)$ group is touched by the algorithm and that each link can reach any point
of the $\Sp(2)$ group, ensuring ergodicity.

We define $A_k$ to be an $\SU(2)$ element belonging to the $k$th $\SU(2)$ subgroup, embedded into $\Sp(2)$. 
For each subgroup in $F$, the $\Sp(2)$ embedding is constructed as follows:

\begin{itemize}
	\item $A_1 = \mqty(
		t_1 & 0 & 0 & t_2 \\
		0 & 1 & 0 & 0 \\
		0 & 0 & 1 & 0 \\
		-t^*_2 & 0 & 0 & t^*_1 
		)$
	\item $A_2 = \mqty(
		1 & 0 & 0 & 0 \\
		0 & t_1 & t_2 & 0 \\
		0 & -t^*_2 & t^*_1 & 0 \\
		0 & 0 & 0 & 1 
		)$
	\item $A_3 = \mqty(
		t_1 & 0 & t_2 & 0 \\
		0 & t_1 & 0 & -t_2 \\
		-t_2^* & 0 & t^*_1 & 0 \\
		0 & t^*_2 & 0 & t^*_1 
		)$
	\item $A_4 = \mqty(
		t_1 & t_2 & 0 & 0 \\
		-t^*_2 & t^*_1 & 0 & 0 \\
		0 & 0 & t_1 & t_2  \\
		0 & 0 & -t^*_2 & t^*_1 
		)$
\end{itemize}

Generating each $A_k$ randomly, we define the new link $U'$ to be

\begin{equation} \label{eq:heatbath_newlink}
	U' = A_4 A_3 A_2 A_1 U,
\end{equation}

since, in this case, the set $F$ contains four $\SU(2)$ subgroups.

If each $A_k$ is randomly distributed as

\begin{equation} \label{eq:probdist}
	P(A_k) = \dd^{(k)}{A_k} \frac{e^{-S[A_k U_{k-1}]}}{\mathcal{Z}_k(U_{k-1})},
\end{equation}

where $\dd^{(k)}$ is the measure of $\SU(2)_k \subset \Sp(2)$, $U_{n} \equiv A_n A_{n-1} \dots A_1 U$ with $U_0 = U$ and 

\begin{equation}
	\mathcal{Z}_k(U) = \int_{\SU(2)_k} \dd{A} e^{-S[AU]} \qc \SU(2)_k \subset \Sp(2),
\end{equation}

then the algorithm will reach thermalization. Since there is no left ideal, we simply need to show that, if $U_{k-1}$
is distributed following a Boltzmann distribution, then the new link $U_k$ will also be distributed in the same way.
Assume that $U_{k-1}$ is distributed as

\begin{equation}
	\dd{P(U_{k-1})} = \frac{1}{\Z} e^{-S[U_{k-1}]}.
\end{equation}

Since the group measure is invariant under multiplication by group elements (which also means 
$\dd^{(k)}{A} = \dd^{(k)}{A^{-1}}$) and

\begin{equation}
	\Z_k(b U) = \Z_k(U) \qc b \in \SU(2)_k,
\end{equation}

then the new link element $U_k = A_k U_{k-1}$ will be distributed as

\begin{equation}
	\dd{P(U_k)} = \int_{\SU(2)_k} \dd^{(k)}{A} \frac{e^{-S[U_k]}}{\Z_k(A^{-1}U_k)} \dd{A^{-1}U_k} \frac{e^{-S[A^{-1}U_k]}}{\Z},
\end{equation}

which is indeed a Boltzmann Distribution. 

At this point, we consider each $a_k$, apply Creutz's algorithm, embed the result back into $\Sp(2)$ thus obtaining
$A_k$, and finally get the new link using \eqref{eq:heatbath_newlink}.

\subsection{Overrelaxation} \label{overrelaxation}

The full algorithm implemented for the simulation also takes advantage of \textit{overrelaxation} \cite{montvay}.
Overrelaxation is used to counter critical slowing down in the vicinity of the phase transition: here,
the system's correlation length is so large that lattice updates of the order of the lattice spacing are
negligible. This causes the simulation to slow down considerably. 
Overrelaxation is a way of choosing new link elements to update the lattice with, which prevents critical
slowing down. 

The idea is to choose a new link element $U'$ "as far as possible" from $U$ (the original link), without actually changing the
action $S$ of the system. \cite{montvay} describes overrelaxation in the case of SU(2) gauge group. For more
complex groups, it is not as straightforward to generate new links. Again, we use Cabibbo and Marinari's idea:
the new link is chosen left multiplying the old link $U$ by randomly generated $A_k$ elements belonging
to Sp(2), which are in turn built embedding SU(2) elements into Sp(2). $k$ labels the SU(2) subgroups and runs
from 1 to 4, as for the heat-bath algorithm. 

Note that overrelaxation is applied at each simulation step along with heat-bath: in fact,
a single complete simulation step consists of one or more steps of overrelaxation and one step of heat-bath, where
a step is the update of the entire lattice exactly once. In our case, overrelaxation is applied three times
before each heat-bath step. This is possible because, as mentioned, overrelaxation doesn't actually change
the action of the system.

Consider $\widetilde{U}_i$ to be one staple surrounding the link $U$, and let us define

\begin{equation}
	R = \sum_{i = 1}^6 \widetilde{U}_i.
\end{equation}

Evidently, there are four staples surrounding each link in a $(2+1)$ dimensional lattice, thus $i = 1,\dots,4$.

\cite{montvay} defines the new link to be
\begin{equation}
	U' = V U^{-1} V,
\end{equation}

where 

\begin{equation}
	V \equiv \det(R)^{\flatfrac{1}{2}} R^{-1},
\end{equation}

the inverse of the projection of $R$ onto the SU(2) gauge group. In the case of the Sp(2) gauge group, 
we define 

\begin{equation} \label{eq:overrelaxation_Ak}
	A_k U_{k-1} = V U_{k-1}^{-1} V,
\end{equation}

instead, where 

\begin{equation}
	U_{n} \equiv A_n A_{n-1} \dots A_1 U \qq{with} U_0 = U.
\end{equation}

Following Cabibbo and Marinari's prescription, the new link will be

\begin{equation} \label{eq:overrelaxation_newlink}
	U' = A_4 A_3 A_2 A_1 U
\end{equation}

in the case of four SU(2) subgroups.

We then extract the SU(2) elements corresponding to the $k$th subgroup of Sp(2) 
from each term in \eqref{eq:overrelaxation_Ak}:

\begin{eqnarray}
	\begin{cases}
		A_k \in \Sp(2) &\longrightarrow a_k \in \SU(2) \\
		U_{k-1} \in \Sp(2) &\longrightarrow u_k \in \SU(2) \\
		\widetilde{U}_i \in \Sp(2) &\longrightarrow r_k^{(i)} \in \SU(2)
	\end{cases}.
\end{eqnarray}

We also define
\begin{equation}
	r_k \equiv \sum_i r_k^{(i)}.
\end{equation}

Using the fact that a sum of SU(2) elements is proportional to a SU(2) element, we can write

\begin{equation}
	u_k r_k = u_k \sum_i r_k^{(i)} = c \bar{u}_k
\end{equation}

with $c = \det(r_k)^{\flatfrac{1}{2}}$ and $\bar{u}_k \in \SU(2)$.

Rewriting \eqref{eq:overrelaxation_Ak} as

\begin{equation}
	A_k = V U^{-1} V U^{-1}
\end{equation}

and applying the equation to the SU(2) elements extracted above, we end up with

\begin{equation}
	\begin{aligned}
		a_k &= \det(r_k) r_k^{-1} u_k^{-1} r_k^{-1} u_k^{-1} \\
		&= \det(r_k) (u_k r_k)^{-1} (u_k r_k)^{-1} \\
		&= \det(r_k) (r_k u_k)^{-2} \\
		&= \det(r_k) (c \bar{u}_k)^{-2}.
	\end{aligned}
\end{equation}

Given that $c = \det(r_k)^{\flatfrac{1}{2}}$, we finally have
\begin{equation}
	a_k = \bar{u}_k^{-2} = \qty(\frac{1}{c} u_k \sum_i r_k^{(i)})^{-2}.
\end{equation}

Exactly as with the heat-bath algorithm, once we find $a_k$, we can embed it into Sp(2) according to the
subgroup it belongs to, which yields the corresponding $A_k$, and find the new link using \eqref{eq:overrelaxation_newlink}.

\newpage
\thispagestyle{empty}
\

\newpage

\section{Results} \label{results}

Here we present finite size scaling analysis and results of our simulations. Lattice simulations were accomplished
by writing a computer program from scratch that implements the Heat-Bath algorithm as presented in section
\ref{heatbath} and the overrelaxation algorithm as presented in section \ref{overrelaxation}. 
The code has been designed to run in parallel on machines with multiple cores in order to
optimize the time taken by the simulation to finish. Simulated lattices have periodic boundary condition in the time
direction in order to define the theory at finite temperature, but this is also the case for the space directions:
this is to simplify calculations for links along the lattice's border and to ignore perimeter effects. This implies
that $t = 0$ and $x_i = 0$ are identified with $t = N_t$ and $x_i = N_s$, respectively.
To update links that are on the edge of the lattice, one must consider links on the other side of the lattice as 
belonging to the staples surrounding the link to update. Another remark is in order: we consider the lattice spacing
for simulations to be $a = 1$. Again, this simplifies simulations and calculations. As discussed at the end of
section \ref{finitetemperature}, setting $a = 1$ means that to change temperature of the system, one has to either
tweak $N_t$, the length of the lattice in the time direction, or $\beta = \flatfrac{8}{g^2}$, the bare gauge
coupling. Ultimately, what we are interested in are temperatures as ratios with respect to $T_c$, the critical 
temperature of the deconfinement transition.

Defining "one step" of either the Heat-Bath or overrelaxation algorithms to be when all links of the lattice
have been update exactly once, our algorithm implements three steps of overrelaxation, to speed up the
decorrelation process and avoid critical slowing down of the simulation, and one step of Heat-Bath, which 
actually changes the action of the system and proceeds with the Monte Carlo simulation. Links are not updated
sequentially, but are updated in parallel to speed up simulations: this has to be done carefully, as this kind
of algorithm is quite susceptible to race conditions. In fact, the algorithm relies on link updates which
depend on the value of the nearest links immediately surrounding it (the staples), but if we change two links
belonging to adjacent lattice sites in parallel, it is not deterministic which one will be updated first. To 
account for this, we have divided lattice sites into even sites and odd sites: given the position 
$x = (t, \va{x})$ of a lattice site in lattice spacing unit, an even (odd) site is one for which the sum of its coordinates
$\sum_{\mu = 0}^d x_\mu$ is an even (odd) integer. If we update all links belonging to all even sites and only
after we take care of odd sites, we are guaranteed to not update concurrently two adjacent sites: by definition,
two adjacent sites $X_1$ and $X_2$ have the same lattice coordinates except one coordinate, in which
the two differ by one lattice spacing:

\begin{equation}
	\begin{aligned}
	X_1 &= (x_0, x_1, \dots, x_i, \dots, x_d, x_{d+1}), \\
	X_2 &= (x_0, x_1, \dots, x_i + 1, \dots, x_d, x_{d+1}).
	\end{aligned}
\end{equation}

This causes the sum of the coordinates to be necessarily even for one site and odd for the other. While this works
perfectly for even values of $N_t$, when $N_t$ is odd a problem presents itself: sites with time coordinates
$t = 0$ and $t = N_t - 1$ and the same spatial coordinates $\va{x}$ are considered 
adjacent due to periodic boundary conditions, but also have the same parity. This problem is solved by first
updating all links which belongs to a single time slice before updating the next time slice. Finally,
we update all links, belonging to even and odd sites, that point in the same direction before moving to the 
next direction. These considerations hold identically for both overrelaxation and Heat-Bath: what changes is the
actual algorithm that updates the single link and not how the link is chosen, or the parallelism is handled.

One final remark is about matrix normalization. Links are represented as $4 \times 4$ matrices which belong
to the $\Sp(2)$ group and must uphold the constraints as discussed in section \ref{sp2}. Due to computer's floating 
point arithmetic, which cannot represent real numbers exactly, operations between real numbers can result in small 
errors, which however can stack up with many sequential operations and cause mistakes in the simulations: multiplying
$\Sp(2)$ matrices many times may cause the final result to not belong to $\Sp(2)$ anymore, despite $\Sp(2)$ being 
closed under multiplication. To account for this, after a certain number of steps, the algorithm
normalizes links matrices to correct possible floating point errors and to make sure they still belong to $\Sp(2)$. 
The algorithm used to normalize matrices is presented in detail in appendix \ref{normalization}.

Each simulation has been made with three steps of overrelaxation followed by one step of Heat-Bath. Matrix normalization
is done every 100 steps of the complete algorithm. The initial configuration of links is arbitrary: often, they
are either the identity element (cold start) or completely randomized (hot start). While this choice
shouldn't make much difference in our case \cite{pepe}, we chose the hot start. To ensure thermalization,
each observable has been measured starting from step 200 of the Monte Carlo history. It has been observed that after 
this number of step, all measured observables set around some average value, indicating that thermalization has been
reached.

Simulated lattices have dimensions for all combinations of $N_t = 5, 6, 7, 8$ and $N_s = 40, 60, 80, 100$: of course,
each of these lattices has dimensions $D = (d+1) = 3$. In order to simulate lattices at different temperatures with
$N_t$ fixed, we have tweaked the value of $\beta$. For each simulation, we have measured the value of the Polyakov
loop by taking the average over the spatial volume at each Monte Carlo step. Figure \ref{fig:mchistory} shows
the Monte Carlo history of the simulation at two different values of $\beta$ for $N_t = 6$, 
corresponding to slightly below and slightly
above the critical $\beta_c$ at which the system has the deconfinement phase transition. It shows two distinct
behaviors: when the system is in the confined phase, the Polyakov loop takes all values around zero, signaling that
the $\mathbb{Z}_2$ center symmetry is unbroken; in the deconfined phase, the Polyakov loop jumps
between two different values from time to time during the Monte Carlo history: this is a sign of the $\mathbb{Z}_2$
center symmetry being broken. 

\begin{figure}[H]
	\centering
	\includegraphics[width=0.7\textwidth]{mchistory_nt=6_l=100_beta=27.png}
	\includegraphics[width=0.7\textwidth]{mchistory_nt=6_l=100_beta=28.png}
	\caption[Monte Carlo history]{Monte Carlo history for the Polyakov loop with $N_t = 6$ and $N_s = 100$ in the
	vicinity of the critical point. Above is a simulation with $\beta = 27$, below with $\beta = 28$}
	\label{fig:mchistory}
\end{figure}

Figure \ref{fig:polyloop_dist} shows the probability distribution of the Polyakov loop $\phi$ before
and after the critical temperature: we can see that the distributions are not the same for different values
of $\beta$. In the
confined phase, the Polyakov loop has a bell-like distribution centered in zero; increasing the value of
$\beta$ (and thus the temperature), the system reaches the critical point: the Polyakov loop distribution starts to
flatten and widen, while being still centered in zero; after the temperature surpasses the critical point, 
the distribution shows two peaks far from zero: this shows the presence of tunneling events for which
the Polyakov loop jumps between two values, which of course mimics the Monte Carlo history in figure \ref{fig:mchistory}.

\begin{figure}[h]
	\centering
	\includegraphics[width=\textwidth]{polydist_nt=6_l=100.png}
	\caption[Polyakov loop distribution]{Distribution of the Polyakov loop for a lattice with $N_t = 6$ and
	$N_s = 100$ at three different values of $\beta$ near the critical value $\beta_c$.}
	\label{fig:polyloop_dist}
\end{figure}

This behavior is typical of second order phase transitions and confirms that for Yang-Mills gauge theory with $\Sp(2)$
gauge group, the deconfinement phase transition is indeed second order, as expected. 

In order to find the precise value of $\beta_c$ for which a deconfinement phase transition occurs, we use
the susceptibility as defined in \eqref{eq:susc_definition}.

This quantity measures the width of the peak (or peaks) of the Polyakov probability distribution. Well in the
confined or deconfined phase, the peak is not very wide. While passing through the transition point, however,
the peak widens: this is where the susceptibility is at the highest value. 

\begin{figure}[h]
	\centering
		\includegraphics[width=0.7\textwidth]{susc_all_l/suscplot_nt=6.png} 

		\includegraphics[width=0.7\textwidth]{susc_all_l/suscplot_nt=8.png}	
	\caption[Susceptibility peaks for $N_t = 6, 8$]{Peak of the susceptibility over spatial volume for
	$N_t = 6$ (above) and $N_t = 8$ (below), for all values of $N_s$ considered. }
	\label{fig:susc_peaks}
\end{figure}

\begin{comment}
\begin{figure}[h]
	\centering
	\begin{tabular}{c c}
		\includegraphics[width=0.45\textwidth]{susc/suscplot_nt=8_l=40.png} &
		\includegraphics[width=0.45\textwidth]{susc/suscplot_nt=8_l=60.png}
		\\
		\includegraphics[width=0.45\textwidth]{susc/suscplot_nt=8_l=80.png} &
		\includegraphics[width=0.45\textwidth]{susc/suscplot_nt=8_l=100.png}
	\end{tabular}	
	\caption[Susceptibility peak for $N_t = 8$]{Peak of the susceptibility over spatial volume for
	$N_t = 8$. From left to right and top to bottom, they are for $N_s = 40, 60, 80, 100$.}
	\label{fig:susc_peaks_nt8}
\end{figure}
\end{comment}

Figure \ref{fig:susc_peaks} shows the susceptibility over spatial volume 
$\flatfrac{\chi}{N_s^2}$ for $N_t = 6, 8$ and all values of $N_s$. To obtain this plot, various
simulations have been made sweeping values of $\beta$ in order to define where the peak is located. We can see
that indeed a peak is present: this signals the occurrence of the deconfinement transition in the vicinity of that
particular value of $\beta$. 

\begin{figure}[h]
	\centering
	\begin{tabular}{c c}
		\includegraphics[width=0.47\textwidth]{susc_fit/suscfit_nt=6_l=40.png} &
		\includegraphics[width=0.47\textwidth]{susc_fit/suscfit_nt=6_l=60.png}
		\\
		\includegraphics[width=0.47\textwidth]{susc_fit/suscfit_nt=6_l=80.png} &
		\includegraphics[width=0.47\textwidth]{susc_fit/suscfit_nt=6_l=100.png}
	\end{tabular}
	\caption[Fit of susceptibility peaks with $N_t = 6$]{Fit of the susceptibility peak for $N_t = 6$ and,
	from left to right and top to bottom, for $N_s = 40, 60, 80, 100$.}
\label{fig:susc_peaks_fit}
\end{figure}

We stress that this analysis has been made for all values
of $N_t = 5, 6, 7, 8$ which produce similar results.

In order to find the precise value of $\beta_c$, we fitted the susceptibility peak
using a model of the form

\begin{equation} \label{eq:quartic_model}
	\beta_c(N_t) \sim a + b(N_t^{(0)} - N_t)^2 + c(N_t^{(0)} - N_t)^3 + d(N_t^{(0)} - N_t)^4,
\end{equation}

where $a,b,c,d,N_t^{(0)}$ represent the parameters of the fit to be found. Figure \ref{fig:susc_peaks_fit} shows
some example of fit of the susceptibility peaks with $N_t = 6$, while table \ref{table:beta_c} lists all values
of $\beta_c$ found for all simulations.

\begin{table}[h]
\begin{center}
	\begin{tabular}{c c}
		\begin{tabular}{|c|c|c|c|}
			\hline
			$N_t$ & $N_s$ & $\beta_c$ & $\chi^2$ \\
			\hline
			\multirow{4}{*}{5} 
			& 40  & $23.312(14)$   & 1.0933\\
			& 60  & $23.2748(52)$ & 1.2741 \\
			& 80  & $23.2886(64)$ & 0.5137 \\
			& 100 & $23.2817(46)$ & 1.4615 \\
			\hline
			\multirow{4}{*}{6} 
			& 40  & $27.589(30)$  & 0.781 \\
			& 60  & $27.547(10)$  & 0.7684 \\
			& 80  & $27.537(12)$ & 0.559 \\
			& 100 & $27.566(13)$ & 1.6649 \\
			\hline
		\end{tabular}
		&
		\begin{tabular}{|c|c|c|c|}
			\hline
			$N_t$ & $N_s$ & $\beta_c$ & $\chi^2$ \\
			\hline
			\multirow{4}{*}{7} 
			& 40  & $32.103(31) $  & 0.1255 \\
			& 60  & $31.8149(92)$ & 0.7616 \\
			& 80  & $31.8190(97)$  & 0.6526 \\
			& 100 & $31.8299(99)$ & 1.328 \\
			\hline
			\multirow{4}{*}{8} 
			& 40  & $36.275(68)$ & 0.7747 \\
			& 60  & $36.103(22)$ & 0.8324 \\
			& 80  & $36.065(19)$ & 1.3907 \\
			& 100 & $36.065(14)$ & 1.0921 \\
			\hline			
		\end{tabular}
	\end{tabular}
\end{center}
\caption[Fitted $\beta_c$]{Value of $\beta_c$ as fitted using \eqref{eq:quartic_model} from the susceptibility peaks.}
\label{table:beta_c}
\end{table}

It is important to take into consideration finite size scaling effects due to volume being finite.
Let us define

\begin{equation}
	x = \frac{\beta}{\beta_c} - 1.
\end{equation}

This variable represents the distance that some simulation defined at $\beta$ has from the critical point. Given
the relation between temperature $T$ and $\beta$, it is also related to the way the observables
defined in section \ref{phase_transition} approach the critical point with their respective critical exponents:

\begin{equation}
	\xi \sim x^{-\nu} \qc \chi \sim x^{-\gamma} \qc \expval{|\phi|} \sim x^{\beta}.
\end{equation}

Finite size scaling provides methods for studying observables in finite volumes and understanding the 
universality class of a statistical system. We define

\begin{equation}
	y = x N_s^{\flatfrac{1}{\nu}} \sim \qty(\frac{N_s}{\xi})^{\flatfrac{1}{\nu}}.
\end{equation}

The correlation length $\xi$ cannot actually be infinite in a finite volume lattice, and it is limited by $N_s$: 
it is thus useful to set a length scale, hence the definition of $y$. We have that the order parameter
$\expval{|\phi|}$ and the susceptibility $N_s^2 \expval{\phi^2}$, if scaled correctly, both follow
two universal scaling curves. Namely, we have

\begin{equation}
	\expval{|\phi|} \sim N_s^{\flatfrac{-\beta}{\nu}} F_1(x N_s^{\flatfrac{1}{\nu}})
\end{equation}

and 

\begin{equation}
	N_s^2 \expval{\phi^2} \sim N_s^{\flatfrac{\gamma}{\nu}} F_2(x N_s^{\flatfrac{1}{\nu}}).
\end{equation}

\begin{figure}[h]
	\centering
	\includegraphics[width=0.7\textwidth]{fss/modphi_fss_nt=6.png}
	\includegraphics[width=0.7\textwidth]{fss/phi2_fss_nt=6.png}
	\caption[Finite size scaling curve for $N_t = 6$]{Universal finite size scaling curves for $N_t = 6$.}
\label{fig:fss}
\end{figure}

\begin{figure}[h]
	\centering
	\includegraphics[width=0.7\textwidth]{adj_fss/adj_modphi_fss_nt=6.png}
	\includegraphics[width=0.7\textwidth]{adj_fss/adj_phi2_fss_nt=6.png}
	\caption[Adjusted finite size scaling curves for $N_t = 6$]{Universal finite size scaling curves for $N_t = 6$,
	with values for $\beta_c$ adjusted in order to lie on the same curve.}
\label{fig:adj_fss}
\end{figure}

\begin{figure}[h]
	\centering
	\begin{tabular}{c c}
		\includegraphics[width=0.45\textwidth]{beta_vs_l/beta_vs_l_nt=5.png} &
		\includegraphics[width=0.45\textwidth]{beta_vs_l/beta_vs_l_nt=6.png}
		\\
		\includegraphics[width=0.45\textwidth]{beta_vs_l/beta_vs_l_nt=7.png} &
		\includegraphics[width=0.45\textwidth]{beta_vs_l/beta_vs_l_nt=8.png}
	\end{tabular}
	\caption[$\beta_c$ as a function of $N_s$]{Values of $\beta_c$ as a function of the spatial length $N_s$: from
	left to right, top to bottom, these are for $N_t = 5, 6, 7, 8$.}
\label{fig:beta_vs_ns}
\end{figure}


Figure \ref{fig:fss} shows the plots of these two curves for $N_t = 6$, where we have set $\nu = 1$, $\gamma = \frac{7}{4}$
and $\beta = \frac{1}{8}$, the critical exponents of the 2D Ising model. If the gauge theory belongs to the same universality
class as the Ising model, these finite size scaling curves $F_1$ and $F_2$ should be universal, and all points should align
perfectly along them. It is clear that for smaller volumes this is not entirely the case: figure \ref{fig:adj_fss} shows
the same curves where $\beta_c$ has been slightly adjusted to fit all points on the same universal curve.
We can see that as $N_s$ gets
bigger, the adjustments are smaller. In fact, figure \ref{fig:beta_vs_ns} shows how
the fitted values for $\beta_c$ change with the spatial length $N_s$: in particular, the values for $N_s = 80$ and
$N_s = 100$ are compatible with each other for all values of $N_t$. 
Finite volume effects are then negligible enough for $N_s = 100$, and we
shall be using results from simulations with spatial volume of $N_s^2 = 100^2$ as a proxy for the thermodynamic limit
observables. 

\begin{figure}[h]
	\centering
	\includegraphics[width=0.7\textwidth]{linear_betafit.png}
	\caption[Linear fit for $\beta_c(N_t)$]{Fit of critical values $\beta_c$ using the linear model 
	$\beta_c(N_t) = a + bN_t$, with $\chi^2 = 2.5612$.}
\label{fig:linear_betafit}
\end{figure}

\begin{figure}[h]
	\centering
	\includegraphics[width=0.7\textwidth]{nonlinear_betafit.png}
	\caption[Non linear fit for $\beta_c(N_t)$]{Fit of critical values $\beta_c$ using the non linear model 
	$\beta_c(N_t) = a + bN_t + \flatfrac{c}{N_t}$, with $\chi^2 = 0.1275$.}
\label{fig:nonlinear_betafit}
\end{figure}

We have fitted the values in table \ref{table:beta_c} for $N_s = 100$ as a function of $N_t$
using two models: one is a simple linear model

\begin{equation} \label{eq:beta_fit}
	\beta_c(N_t) \sim a + bN_t \qq{with}
	\begin{cases}
		a = 1.947(22) \\
		b = 4.2674(38)
	\end{cases},
\end{equation}

shown in figure \ref{fig:linear_betafit}, while the other model is a non linear model:

\begin{equation}
	\beta_c(N_t) \sim a + b N_t + \frac{c}{N_t} \qq{with} 
	\begin{cases}
		a = 2.99(47) \\
		b = 4.183(38) \\
		c = -3.1 \pm 1.4
	\end{cases}.
\end{equation}

This comes from \cite{beta-fit} and is shown in figure \ref{fig:nonlinear_betafit}. Given that the $\chi^2$ test
for the linear model yields $\chi^2 = 2.5612$ and $\chi^2 = 0.1275$ for the nonlinear model, and given the large error
on the $c$ coefficient, we declare our data insufficient to accurately reproduce the non linear deviation of $\beta_c(N_t)$,
and proceed to use the linear fit from now on.

\begin{table}[h]
	\centering
	\begin{tabular}{|c|c|c|c|c|c|}
		\hline
		$N_t$ & $N_s$ & $\flatfrac{T}{T_c} $ & $R$ & $\xi$ & $\chi^2$ \\
		\hline
		\multirow{6}{*}{6} 
			& \multirow{2}{*}{80}
				& \multirow{2}{*}{0.95}
					& $(2, 14)$  & $18.106(92)$ & 1.5862 \\
				& & & $(16, 38)$ & $18.09(12)$  & 0.6533 \\
			\cline{2-6}
			& \multirow{4}{*}{100} 
				& \multirow{2}{*}{0.86} 
					& $(2, 6)$   & $7.422(24)$ & 3.8806 \\
					& & & $(8, 50)$ & $7.861(18)$  & 1.1264 \\
					\cline{3-6}
				& & \multirow{2}{*}{0.89} 
					& $(3, 8)$   & $9.359(51)$ & 2.3199 \\
					& & & $(8, 25)$  & $9.752(28)$ & 2.3932 \\
		\hline
		\multirow{4}{*}{8} & \multirow{4}{*}{100}
				& \multirow{2}{*}{0.86}
					& $(3, 7)$  & $10.543(65)$ & 0.5551 \\
				& & & $(10, 20)$ & $10.615(49)$  & 2.4229 \\
			\cline{3-6}
				& & \multirow{2}{*}{0.89}
					& $(3, 7)$   & $13.17(12)$  & 0.6701 \\
				& & & $(11, 35)$ & $13.819(52)$ & 1.3252 \\
		\hline
		\multirow{2}{*}{5} & \multirow{2}{*}{100} & \multirow{2}{*}{0.95}
			    & $(2, 13)$ & $14.157(66)$ & 1.433 \\
			& & & $(12, 50)$ & $14.26(58)$ & 1.4986 \\
		\hline
		\multirow{2}{*}{6} & \multirow{2}{*}{100} & \multirow{2}{*}{0.95}
			    & $(2, 17)$ & $17.42(11)$ & 0.8582 \\
			& & & $(14, 42)$ & $16.86(12)$ & 0.513 \\
		\hline
		\multirow{2}{*}{7} & \multirow{2}{*}{100} & \multirow{2}{*}{0.95}
			    & $(2, 20)$ & $20.97(15)$ & 0.5012 \\
			& & & $(17, 50)$ & $20.66(17)$ & 0.4379 \\
		\hline
		\multirow{2}{*}{8} & \multirow{2}{*}{100} & \multirow{2}{*}{0.95}
			    & $(1, 24)$ & $24.95(22)$ & 2.9113 \\
			& & & $(18, 50)$ & $24.26(28)$ & 0.4311 \\
		\hline
	\end{tabular}
\caption[Fitted correlation length]{Values for the correlation length $\xi$ found by fitting the short and large distance
	Polyakov loop two point function with the spin-spin correlator.}
\label{table:corr_length}
\end{table}
\begin{comment}

\begin{table}[h]
	\centering
	\begin{tabular}{|c|c|c|c|c|c|}
		\hline
		$N_t$ & $N_s$ & $\flatfrac{T}{T_c} $ & $R$ & $\xi$ & $\chi^2$ \\
		\hline
		\multirow{6}{*}{6} 
			& \multirow{2}{*}{80}
				& \multirow{2}{*}{0.95}
					& $(2, 14)$  & $18.106(92)$ & 1.5862 \\
				& & & $(16, 38)$ & $18.09(12)$  & 0.6533 \\
			\cline{2-6}
			& \multirow{4}{*}{100} 
				& \multirow{2}{*}{0.86} 
					& $(2, 6)$   & $7.422(24)$ & 3.8806 \\
					& & & $(8, 50)$ & $7.861(18)$  & 1.1264 \\
					\cline{3-6}
				& & \multirow{2}{*}{0.89} 
					& $(3, 8)$   & $9.359(51)$ & 2.3199 \\
					& & & $(8, 25)$  & $9.438(27)$ & 8.3842 \\
		\hline
		\multirow{4}{*}{8} & \multirow{4}{*}{100}
				& \multirow{2}{*}{0.86}
					& $(3, 7)$  & $10.543(65)$ & 0.5551 \\
				& & & $(10, 20)$ & $10.615(49)$  & 2.4229 \\
			\cline{3-6}
				& & \multirow{2}{*}{0.89}
					& $(3, 7)$   & $13.17(12)$  & 0.6701 \\
				& & & $(11, 35)$ & $13.265(48)$ & 4.2232 \\
		\hline
		\multirow{2}{*}{5} & \multirow{2}{*}{100} & \multirow{2}{*}{0.95}
			    & $(2, 13)$ & $14.157(66)$ & 1.433 \\
			& & & $(14, 32)$ & $14.31(10)$ & 0.3525 \\
		\hline
		\multirow{2}{*}{6} & \multirow{2}{*}{100} & \multirow{2}{*}{0.95}
			    & $(2, 17)$ & $17.42(11)$ & 0.8582 \\
			& & & $(18, 37)$ & $17.28(22)$ & 0.0261 \\
		\hline
		\multirow{2}{*}{7} & \multirow{2}{*}{100} & \multirow{2}{*}{0.95}
			    & $(2, 20)$ & $20.97(15)$ & 0.5012 \\
			& & & $(21, 47)$ & $21.22(28)$ & 0.0661 \\
		\hline
		\multirow{2}{*}{8} & \multirow{2}{*}{100} & \multirow{2}{*}{0.95}
			    & $(1, 24)$ & $24.95(22)$ & 2.9113 \\
			& & & $(25, 47)$ & $25.78(72)$ & 0.0242 \\
		\hline
	\end{tabular}
\caption[Fitted correlation length]{Values for the correlation length $\xi$ found by fitting the short and large distance
	Polyakov loop two point function with the spin-spin correlator.}
\label{table:corr_length}
\end{table}
\end{comment}

\eqref{eq:beta_fit} represents the critical value of the coupling $\beta$ at which the deconfinement transition occurs.
Using the way we have defined the critical temperature $T_c$ in \eqref{eq:critical_temperature}, we can invert $\beta_c(N_t)$
to find the critical value $N_{t,c}(\beta)$: if we take its reciprocal, we find $T_c$ for a fixed value of $N_t$ (recall
that we have set the lattice spacing $a = 1$). Now we can do simulations for $N_t = 5, 6, 7, 8$ and $N_s = 100$ for a value of
$\beta$ that, for each choice of $N_t$, provides the gauge theory at temperature slightly below $T_c$ at which we can study
the two point function of the Polyakov loop defined in \eqref{eq:polyloop_correlator}. In the path integral formulation,
this is

\begin{equation} \label{eq:polyloop_correlator_pathintegral}
	\expval{\phi(0) \phi(R)} = \frac{1}{\Z} \int \D U e^{-S_W[U]} \frac{1}{N_s^2} \sum_{\va{x}, \va{y}}\phi(\va{x}) \phi(\va{y}),
\end{equation}

where $|\va{x} - \va{y}| = R$.

\begin{figure}[h]
	\centering
	\includegraphics[width=0.7\textwidth]{corr_t095/corr_short_fit_nt6.png}
	\includegraphics[width=0.7\textwidth]{corr_t095/corr_long_fit_nt6.png}
	\caption[Polyakov loop correlator fit for $N_t=6$]{
		Fit of the Polyakov loop two point function using the spin-spin correlator from the 2D Ising model
		for $N_t = 6$ and $N_s = 100$. Above is the short distance $\xi < R$ fit; below is the large distance $\xi > R$ fit.
	}
\label{fig:corr_fit_nt6}
\end{figure}

\begin{figure}[h]
	\centering
	\includegraphics[width=0.7\textwidth]{corr_t095/corr_short_fit_nt8.png}
	\includegraphics[width=0.7\textwidth]{corr_t095/corr_long_fit_nt8.png}
	\caption[Polyakov loop correlator fit for $N_t=8$]{
		Fit of the Polyakov loop two point function using the spin-spin correlator from the 2D Ising model
		for $N_t = 8$ and $N_s = 100$. Above is the short distance $\xi < R$ fit; below is the large distance $\xi > R$ fit.
	}
\label{fig:corr_fit_nt8}
\end{figure}


If the Svetisky and Yaffe conjecture is true, and the system belongs to the universality class of the 2D Ising model, the 
short distance ($R < \xi$) and large distance ($R > \xi$) behavior of \eqref{eq:polyloop_correlator_pathintegral} 
in the vicinity of the critical point should be the same of the spin-spin
correlator in \eqref{eq:short_distance_spin} and \eqref{eq:large_distance_spin}. Table \ref{table:corr_length} lists the fitted
values of $\xi$ at different temperatures and values of $N_t$; figures \ref{fig:corr_fit_nt6} and
\ref{fig:corr_fit_nt8} show the short and large distance fit for $N_t = 6$ and $N_t = 8$, respectively.

There is some degree of agreement in our data, and correlations lengths tend to be compatible in the short
and large distance behavior. This confirms that the Svetisky and Yaffe conjecture
is a good description of the theory near the critical point of the deconfinement phase transition, because the theory
is in fact in the universality class of the 2D Ising model. 
Also, these results are in line with the EST description
of the gauge theory, which predicted the large distance behavior in \eqref{eq:est_polyloop_corr} for the two point 
correlation function of Polyakov loops.

We can also see that the measurement is not easy: when $\xi$ is low, there
are not many values of $R$ to choose for the fit in the short distance regime. $R$ can take only integer values on the lattice,
and when $\xi$ is small, changing $R$ by only 1 can drastically change the fit. In the large distance regime, instead, the
uncertainty of the two point function tends to grow, and we would require more iterations of the simulations in order to reduce
the error associated to $\xi$ and find a better fit. In fact, for $N_t = 6$ and $N_s = 80$, simulations are cheaper and quicker,
enabling us to do more iterations and collect more data, thus yielding a better fit. As it was the case for the short distance
regime, $R$ can be quite limited in the large distance regime, also. In fact, for $N_t = 8$, $N_s = 100$ and $T = 0.95T_c$,
the fit is heavily limited by the large value of $\xi$: since we are using periodic boundary conditions and taking into 
account the mirror effect by adding a term to the spin-spin correlator calculated
in $N_s - R$ (as shown in \eqref{eq:large_distance_spin}), for a lattice with $N_s = 100$, the maximum distance $R$ the
correlator can reach is $R = \flatfrac{N_s}{2} = 50$. For $\xi \sim 25$, this limits the range that $R$ can span in order
to fit the correlator. Choosing higher values of $N_s$, while keeping $N_t$ small, can improve the fit by giving more room
for adjustments of the distance $R$: of course, this would translate to more complex simulations, given that the lattice
volume increases. 

\FloatBarrier

\newpage
\thispagestyle{empty}
\

\newpage

\section{Conclusions} \label{conclusions}

We have described how a Yang-Mills gauge theory can be simulated on a lattice, starting from theoretical principles
that enables us to make the theory real and well defined and treatable with statistical methods, i.e. Monte Carlo methods.
We have also shown how importance sampling can be used to measure observables in the context of the lattice regularization
of the gauge theory, and how to actually generate lattice configurations that respect the Boltzmann probability 
distribution using Markov chains.

We have also described in detail the Heat-Bath algorithm, a very efficient way to simulate a LGT for a generic gauge
group which can be decomposed into $\SU(2)$ subgroups. The gauge group of choice for this work has been $\Sp(2)$ due to the 
fact that its center is $\mathbb{Z}_2$ regardless of the group size. Overrelaxation was also taken into account, described
and implemented: this further enhances the speed of the algorithm by speeding up decorrelation and avoiding critical 
slowing down in the vicinity of the phase transition. 

We successfully implemented the Heat-Bath and overrelaxation algorithms for the $(2+1)$ dimensional
Yang-Mills gauge theory regularized on the lattice: our simulations show that the LGT we have simulated
does in fact exhibit a second order deconfinement phase transition, as expected for the $\Sp(2)$ group and predicted by
the Svetisky and Yaffe conjecture. We have tested the conjecture by measuring the critical temperature at different lattice
sizes and simulating the LGT for temperature values that are slightly below the critical temperature. Data shows that
indeed the (2+1) dimensional Yang-Mills at finite temperature with gauge group $\Sp(2)$ is in the same universality class
of the 2D Ising model. In the vicinity of the critical point, the spin-spin correlator describes well the behavior
of the Polyakov loop two point function: thus, the mapping by Svetisky and Yaffe described in section \ref{conjecture} appears
to be valid, even in the case of the $\Sp(2)$ gauge group. Also, the large distance behavior of the Polyakov loop correlator
is predicted by the EST description at large distaces: this is in line with the idea that EST is a good effective description
of the theory far from the critical point, in the large distance limit. Near the critical point, the flux tube
is disturbed by critical effects and the model of a vibrating string becomes a bad approximation.

Going forward, it is quite interesting to increase the data collected in terms of iterations number but also increase 
the spatial volume of the lattice: this would make the measurement of the two point function of the Polyakov loop more
adjustable and precise and would yield a more accurate calculation for the correlation length. Changing the 
gauge group to $\Sp(N)$ with $N > 2$ would also be very interesting: the center symmetry would be the same, and keeping the
lattice dimensions at $(2+1)$ would imply that, according to the Svetisky-Yaffe conjecture, the system would still
be in the 2D Ising model universality class while also having a bigger gauge group. This would be an amazing test for the
conjecture itself. 

Also, the Polyakov loop correlator can be used to study the interquark potential: it would be interesting to accurately measure
this quantity in the confined phase, but with high temperature, to push the limits of the EST description. 
In particular, more work could be done to search for correction
terms to the interquark potential (like the L{\"u}scher term) and to the EST model in general. We have seen that the
Nambu-Goto model describes the theory quite well at large distances while not being as reliable in the vicinity
of the critical point: the Polyakov loop correlator and its relationship with the interquark potential could be used
in order to find corrections terms and to see how the system gets from being well described by the EST to being well
described by the Svetisky and Yaffe conjecture near the critical point.
Also, the EST makes prediction about $\sigma(T)$, the string tension
above for temperatures above $T = 0$, which appears in the interquark potential: 
this is also an area that could be more interesting to look into using Monte
Carlo simulations in the case of the $\Sp(2)$ gauge group.


\newpage

\begin{appendices}
\section{Matrix normalization} \label{normalization}
	
Here we describe in detail how matrices are normalized during simulations. As mentioned in section \ref{results},
computer arithmetic can lead to small errors due to floating numbers representation. If matrices are not renormalized
from time to time, they risk deviating from the $\Sp(2)$ group to which they should belong. 

An $\Sp(2)$ matrix is also an $\SU(4)$ matrix: this means that a link $U$ must be unitary and have determinant 1:

\begin{equation}
	UU^\dagger = \id \qc \det U = 1.
\end{equation}	

Now, consider the $\Sp(2)$ matrix in the form \eqref{eq:sp2_matrix_form}
 and let us define the 1st, 2nd and 4th rows of this matrix as

\begin{equation}	
	\begin{aligned}
		\va{V}_1 &= \mqty(W_{11} & W_{12} & X_{11} & X_{12}) \\
		\va{V}_2 &= \mqty(W_{21} & W_{22} & X_{21} & X_{22}) \\
		\va{V}_3 &= \mqty(-X_{12}^* & X_{11}^* & -W_{12}^* & W_{11}^*).
	\end{aligned}
\end{equation}

Calculating $UU^\dagger$, we get
\begin{equation}
	UU^\dagger = 
	\mqty(
		\abs{\va{V}_1}^2 & \va{V}_2^*\cdot\va{V}_1 & \qty(\va{V}_2^*\cdot\va{V}_3)^* & 0 \\
		\qty(\va{V}_2^*\cdot\va{V}_1)^* & \abs{\va{V}_2}^2 & 0 & \qty(\va{V}_2^*\cdot\va{V}_3)^* \\
		\va{V}_2^*\cdot\va{V}_3 & 0 & \abs{\va{V}_2}^2 & -\va{V}_2^*\cdot\va{V}_1 \\
		0 & \va{V}_2^*\cdot\va{V}_3 & -\qty(\va{V}_2^*\cdot\va{V}_1)^* & \abs{\va{V}_1}^2
	)
\end{equation}

noting that $\abs{\va{V}_1}^2 = \abs{\va{V}_3}^2$.

We want this matrix to be the identity. The solution is to scale $\va{V}_1$ and $\va{V}_3$ by their norm, $\abs{\va{V}_1}$:
\begin{equation}
\begin{aligned}
\va{V}_1 &\longrightarrow \va{V'}_1 = \frac{\va{V}_1}{\abs{\va{V}_1}} \\
\va{V}_3 &\longrightarrow \va{V'}_3 = \frac{\va{V}_3}{\abs{\va{V}_3}} = \frac{\va{V}_3}{\abs{\va{V}_1}}.
\end{aligned}
\end{equation}

Now the first and the last elements are equal 1. Then, we substitute $\va{V}_2$ with $\va{V'}_2$ defined as

\begin{equation}
\va{V}_2 \longrightarrow \va{V'}_2 = \va{V}_2 - \qty(\va{V}_2^*\cdot\va{V'}_1)^*\va{V'}_1 - \qty(\va{V}_2^*\cdot\va{V'}_3)^*\va{V'}_3.
\end{equation}

Considering that $\va{V'}_3^*\cdot\va{V'}_1 = 0$ and that now $\abs{\va{V'}_1} = \abs{\va{V'}_3} = 1$, we can check that

\begin{equation}
\begin{aligned}
\va{V'}_2^* \cdot \va{V'}_1 &= \va{V}_2^* \cdot \va{V'}_1 - \qty(\va{V}_2^* \cdot \va{V'}_1)\qty(\va{V'}_1^*\cdot\va{V'}_1) - \qty(\va{V}_2^* \cdot \va{V'}_3)\qty(\va{V'}_3^*\cdot\va{V'}_1) \\
&= \va{V}_2^* \cdot \va{V'}_1 - \qty(\va{V}_2^* \cdot \va{V'}_1)\abs{\va{V'}_1}^2 \\
&= \va{V}_2^* \cdot \va{V'}_1 - \va{V}_2^* \cdot \va{V'}_1 \\ 
&= 0.
\end{aligned}
\end{equation}

The same thing happens when multiplying by $\va{V'}_3$: 

\begin{equation}
\va{V'}_2^* \cdot \va{V'}_3 = 0.
\end{equation}

Finally, we can simply scale $\va{V'}_2$ by its norm, so that the matrix diagonal is all 1s:

\begin{equation}
\va{V'}_2 \longrightarrow \va{V_2''} = \frac{\va{V'}_2}{\abs{\va{V'}_2}},
\end{equation}

and rebuild the original $U$ matrix using the $2\times2$ matrices $W$ and $X$:
\begin{equation}
W = \mqty(V_1'^{(1)} & V_1'^{(2)} \\ V_2''^{(1)} & V_2''^{(2)}) \qc
X = \mqty(V_1'^{(3)} & V_1'^{(4)} \\ V_2''^{(3)} & V_2''^{(4)}).
\end{equation}

\end{appendices}

\newpage
\thispagestyle{empty}
\

\newpage

\printbibliography
% https://www.overleaf.com/learn/latex/Bibliography_management_in_LaTeX

\end{document}